\ncertchapter{10}{Circular Dynamics and Non-Inertial Frames}\chlabel{ch:circdyn}

\chref{ch:circkin} proved that a body moving on a circle must have an inward acceleration $v^2/r$ (speed $v$, radius $r$). It said nothing
about what produces that acceleration. This chapter asks, for each real situation, \emph{which actual
force supplies it}? For a car on a flat road it is friction; for a conical pendulum, a component of the
string's tension; for a satellite, gravity; for a rider in a rotor, the normal force from the wall. The
physics is the same in every case.

The chapter works through horizontal circles, vehicles on curved roads, and the vertical circle, which is
harder because the speed varies around the loop and a string can go slack. It then describes motion from
inside an accelerating or rotating frame, where Newton's laws fail unless a pseudo-force is added, and
closes with an application that draws these threads together: how a spinning spacecraft produces
artificial gravity, and why any practical design must be very large.

\section{The Central Principle}\label{sec:centralprinciple}

A beginner's instinct is to draw a free-body diagram and then add ``the centripetal force'' as an extra
arrow. \textbf{That is always wrong.} There is no such force in nature, and the extra arrow counts a
force twice, an error that is hard to find later because each individual step looks reasonable. What
\chref{ch:circkin} established is a \emph{requirement}:
\[
\text{if a body of mass } m \text{ moves on a circle of radius } r \text{ at speed } v,
\quad\text{then}\quad
\sum F_{\text{radial}} \;=\; \frac{mv^{2}}{r}.
\]
The left side is a sum of \emph{real} forces (tension, gravity, friction, the normal force). The right
side is the bill those forces must collectively settle. ``Centripetal'' names the \emph{role}, like
``goalkeeper''; it never names a new player.

Two beliefs about a stone whirled on a string are worth correcting now. The outward tug you feel in your
hand is real, but it is the third-law partner of the string's inward pull on the stone: it acts on
\emph{you}, not on the stone. And if the string snaps, the stone is not flung radially outward. Nothing
was pushing it outward; the inward pull simply stops, and by the first law (\chref{ch:newton}) the stone
carries on in a straight line along the \emph{tangent} at the point of release.

So the method is always the same two steps: \textbf{(1)} list the real forces; \textbf{(2)} resolve them
along the radius and set the inward total equal to $mv^{2}/r$. If a tangential acceleration is also
present, resolve along the tangent too and set that total to $ma_t$. Throughout the chapter these
equations are written in the ground frame, which we treat as inertial; pseudo-forces appear only in
\S\ref{sec:noninertial}, where the frame itself accelerates.

\begin{methodbox}{The working equations}
\[
\sum F_{\text{radial (inward)}}=\frac{mv^{2}}{r}=m\omega^{2}r,
\qquad\qquad
\sum F_{\text{tangential}}=ma_t=mr\alpha.
\]
Here $\omega=v/r$ is the angular speed (\si{rad/s}), $a_t$ the tangential acceleration and $\alpha$ the
angular acceleration (\si{rad/s^2}), all as in \chref{ch:circkin}. In \emph{uniform} circular motion the
second equation is zero, and only the radial one is needed.
\end{methodbox}

\section{Horizontal Circles}
In every case of this section the body stays at one height, so the vertical forces balance and the
horizontal ones supply $mv^{2}/r$; the cases differ only in which real force has the horizontal part.

\subsection{The conical pendulum: how can it just stay there?}

A bob of mass $m$ on a string of length $\ell$ is whirled so that it travels a horizontal circle while the
string, with tension $T$, sweeps out a cone at a fixed angle $\theta$ to the vertical (Figure~\ref{fig:conical}). Students find this puzzling.

\begin{motivbox}{Resolving the puzzle: it is \emph{not} in equilibrium}
Three questions are usually tangled together.

\smallskip
\textbf{``Why does it not fall?''} Because vertically it \emph{is} in equilibrium. The bob never moves
up or down, so its vertical acceleration is zero and the vertical forces cancel: the upward part of the
tension, $T\cos\theta$, matches the weight $mg$.

\smallskip
\textbf{``Then why does it not swing back to the middle?''} Because horizontally it is \emph{not} in
equilibrium, and is not supposed to be. The inward force $T\sin\theta$ is unopposed, and that is
exactly right: the bob \emph{is} accelerating, continuously bent out of the straight line it would
otherwise follow. The unbalanced inward force is the reason the bob goes round instead of off.

\smallskip
\textbf{``Why does it sit at that particular angle?''} The angle is \emph{selected} by the speed. Spin
the bob faster and it needs a bigger inward force; it swings outward, which increases $\theta$ and with
it $T\sin\theta$, until supply matches demand again. Slow it down and it falls inward to a smaller cone.

\smallskip
\emph{In short:} \textbf{vertical equilibrium, horizontal acceleration}, as in every horizontal circle
of this section.
\end{motivbox}

\begin{figure}[htbp]
\centering
\tikzfig{cd_conical}
\figcaption{fig:conical}{The conical pendulum. Only two real forces act: the tension $T$ along the
string and the weight $mg$. Resolving $T$ vertically supports the bob; resolving it horizontally bends
the bob's path. The circle has radius $r=\ell\sin\theta$, \emph{not} $\ell$.}
\end{figure}

\begin{derivbox}{Speed, angular speed and period}
\stepn{1} \textbf{Identify the real forces.} Exactly two: the tension $T$ along the string, and the
weight $mg$ downward. (No ``centripetal force'' is drawn; see \S\ref{sec:centralprinciple}.)

\stepn{2} \textbf{Vertical: equilibrium.} The bob stays at the same height, so
\[
T\cos\theta=mg. \tag{i}
\]

\stepn{3} \textbf{Horizontal: Newton's second law along the radius.} The only horizontal force is the
inward component of the tension, and the radius of the circle is $r=\ell\sin\theta$:
\[
T\sin\theta=\frac{mv^{2}}{r}. \tag{ii}
\]

\stepn{4} \textbf{Divide (ii) by (i)}: this eliminates both $T$ and $m$ in one stroke:
\[
\tan\theta=\frac{v^{2}}{rg}
\qquad\Longrightarrow\qquad
\boxed{\;v=\sqrt{rg\tan\theta}\;}
\]

\stepn{5} \textbf{Angular speed and period.} With $\omega=v/r$, and writing the period as $T_{\text{per}}$
to keep it distinct from the tension $T$,
\[
\omega=\sqrt{\frac{g\tan\theta}{r}},
\qquad
T_{\text{per}}=\frac{2\pi}{\omega}=2\pi\sqrt{\frac{r}{g\tan\theta}}.
\]
Substituting $r=\ell\sin\theta$ and simplifying $\dfrac{\sin\theta}{\tan\theta}=\cos\theta$:
\[
\boxed{\;T_{\text{per}}=2\pi\sqrt{\frac{\ell\cos\theta}{g}}\;}
\]

\stepn{6} \textbf{Sanity checks.}
\begin{itemize}
\item As $\theta\to0$ (a nearly vertical string, very slow whirl), $T_{\text{per}}\to2\pi\sqrt{\ell/g}$,
      the period of a simple pendulum swinging with \emph{small} amplitude (derived when oscillations
      are studied). The match is no accident: a small conical motion is two small swings at right
      angles, a quarter-cycle out of step. \checkmark
\item As $\theta\to90^\circ$, $\cos\theta\to0$ so $T_{\text{per}}\to0$ and, from (i),
      $T=mg/\cos\theta\to\infty$. A perfectly horizontal string is impossible: it has no upward
      component to hold the bob up, however fast the bob is spun. \checkmark
\item The mass has cancelled: on strings of the same length, a heavy bob and a light one cone out at
      the same angle for the same speed. \checkmark
\end{itemize}
\end{derivbox}

\begin{trapbox}{the radius is $\ell\sin\theta$, not $\ell$}
The string length is the \emph{slant}; the circle's radius is its horizontal projection. Using $r=\ell$
is the most common error in this problem.
\end{trapbox}

\subsection{A particle inside a smooth cone}
In the conical pendulum the slanted force was a tension. A surface can play the same part. Picture a funnel held point-down, with a marble running round the inside of it. Given the right speed
the marble settles into a steady horizontal circle partway up the wall, neither climbing nor slipping
down. The apparatus looks nothing like a conical pendulum, but the force diagram is the same one
(Figure~\ref{fig:cdcone}).

\begin{figure}[htbp]
\centering
\tikzfig{cd_cone}
\figcaption{fig:cdcone}{A particle riding the inside of a smooth cone whose wall makes an angle $\theta$ with the horizontal. The normal force is perpendicular to the wall, so it makes $\theta$ with the vertical: its vertical part supports the particle and its horizontal part turns it, exactly as the tension did in the conical pendulum.}
\end{figure}

\begin{derivbox}{The same two steps, a different force}
A particle rides a horizontal circle of radius $r$ on the smooth inside wall of a cone whose axis is
vertical and whose wall makes an angle $\theta$ with the horizontal (so the normal force makes
$\theta$ with the vertical, and the cone's half-angle is $90^\circ-\theta$). The wall is smooth, so
again only \emph{two} real forces act: the normal force $N$ (perpendicular to the wall) and the
weight $mg$.

\stepn{1} \textbf{Vertical equilibrium:} $N\cos\theta=mg$.

\stepn{2} \textbf{Radial (inward) acceleration:} $N\sin\theta=\dfrac{mv^{2}}{r}$.

\stepn{3} \textbf{Divide:} $\tan\theta=\dfrac{v^{2}}{rg}$, so
\[
\boxed{\;v=\sqrt{rg\tan\theta}\;}
\]
This is \emph{identical} to the conical pendulum.

\smallskip
\textbf{Watch which angle is given.} Many problems specify the cone by its half-angle $\alpha$, measured
from the axis. Then $\theta=90^\circ-\alpha$ and the same result reads $v=\sqrt{rg\cot\alpha}$.

\smallskip
\textbf{Why identical?} One slanted force does double duty, its vertical part holding the body up
and its horizontal part turning it; whether it is a tension or a normal force does not matter.
\end{derivbox}

\subsection{The rotor and the death well}\label{sub:rotor}
So far one slanted force both held the body up and turned it. Remove the slant, so that the surface is
vertical, and the two jobs must be shared between two different forces. Two fairground machines pose this problem.
\textbf{The death well} (the \emph{maut ka kuan} of travelling fairs) is a large wooden cylinder
standing on its end so that its wall is vertical; a motorcyclist rides round the
\emph{inside} of that wall, high above the floor. Here the \textbf{wall is stationary and the rider
moves}. \textbf{The rotor} uses the same cylinder, but the \textbf{cylinder itself
spins} about its vertical axis while the riders stand with their backs against the wall. Once it is
turning fast enough the floor is lowered away, and the riders stay pinned to the wall.

In both, a body of mass $m$ travels a horizontal circle of radius $r$ pressed against a \emph{vertical}
surface, and the question is the same: \emph{how fast must the motion be for the body not to slide down
the wall?} What makes the problem feel unfamiliar is a reversal of geometry. On a floor the normal force
is vertical and friction is horizontal; on a vertical wall \textbf{the normal force is horizontal and
friction is vertical} (Figure~\ref{fig:rotor}).

\begin{figure}[htbp]
\centering
\tikzfig{cd_rotor}
\figcaption{fig:rotor}{Two fairground rides that look different and are not. In the \emph{death well} the
wall is fixed and the rider drives round the inside; in the \emph{rotor} the cylinder spins and carries
the rider with it. In both, the normal force is horizontal and friction vertical.}
\end{figure}

\begin{derivbox}{The minimum speed not to slide down}
\stepn{1} \textbf{Get the directions right.} The wall is vertical, so the normal force $N$ it
exerts is \textbf{horizontal}, pointing inward at the axis.
Friction acts along the surface, which here means \textbf{vertically}, and it must point \emph{up} to
stop the rider sliding down.

\stepn{2} \textbf{Radial equation.} The only horizontal force is $N$, and it does the turning:
\[
N=\frac{mv^{2}}{r}=m\omega^{2}r.
\]

\stepn{3} \textbf{Vertical equation.} The rider does not slide, so friction balances the weight exactly:
\[
f=mg.
\]

\stepn{4} \textbf{Impose the friction limit} $f\le\mu N$ (here $\mu$ is the static coefficient, since
the rider must not start to slip):
\[
mg\le\mu\,m\omega^{2}r
\qquad\Longrightarrow\qquad
\boxed{\;\omega_{\min}=\sqrt{\frac{g}{\mu r}}\;}
\]
The mass cancels: a heavy rider and a light one are equally safe, because both the weight to be
supported and the normal force (hence the friction available) scale with $m$.

\smallskip
\textbf{The counter-intuitive reading:} spinning \emph{faster} presses you harder against the wall,
which buys more friction, which is what holds you up. The ride's safety comes from its speed.
\end{derivbox}

\subsection{Why a cyclist leans into the bend}
In the rotor friction held the body up; on a level road it does the turning, and that raises a new
question. A cyclist taking a corner at speed tips the whole machine inward; a rider who tried to stay upright
would fall \emph{outward}. Yet the inward force comes from friction at the tyres, which is available
whether the bicycle leans or not. The reason for leaning involves \emph{torque}, the turning effect of a
force, rather than force itself: it is about not toppling while using the inward force.

\begin{figure}[htbp]
\centering
\tikzfig{cd_cyclist}
\figcaption{fig:cdcyc}{Why leaning is compulsory, not a matter of style (seen from behind, centre of the bend to the left). Upright (left), friction $f$ acts at road level, a height $h$ below $G$, and its moment about $G$ topples the rider outwards. Leaning (right) by $\theta$ with $\tan\theta=f/N$ makes the resultant contact force $F$ (dashed: $N$ and $f$ added) pass through $G$, so it has no moment about $G$.}
\end{figure}

\begin{derivbox}{The lean angle}
Treat the rider and bicycle as one rigid body of mass $m$ moving at steady speed $v$ round a level bend
of radius $r$. Friction acts at the \emph{ground}, well below the centre of gravity, so it would topple
the rider outward unless something compensates. Leaning is that compensation (Figure~\ref{fig:cdcyc}).

\stepn{1} \textbf{The two forces at the ground} are the normal force $N$ (up) and friction $f$
(inward). Their resultant must pass through the centre of gravity $G$, the point at which the weight
acts (in uniform gravity, the centre of mass of \chref{ch:com}); otherwise there is a net torque
about $G$ and the cyclist starts to rotate, that is, falls over. (Torque is treated fully in
\chref{ch:rotation}. Here we need one fact: a body that is not starting to rotate has zero net torque,
or \emph{moment}, about $G$, even while $G$ accelerates. About any other point this fails for an
accelerating body.)

\stepn{2} \textbf{Vertical:} $N=mg$. \quad \textbf{Radial:} $f=\dfrac{mv^{2}}{r}$.

\stepn{3} \textbf{The lean angle} $\theta$ from the vertical is the angle of that resultant:
\[
\tan\theta=\frac{f}{N}=\frac{mv^{2}/r}{mg}
\qquad\Longrightarrow\qquad
\boxed{\;\tan\theta=\frac{v^{2}}{rg}\;}
\]

\stepn{4} \textbf{The limit.} Since $f\le\mu N=\mu mg$, $\tan\theta\le\mu$: the steepest possible lean is
$\tan^{-1}\mu$, and a bend that demands more lean than that cannot be taken at that speed however
skilled the rider.

\smallskip
\textbf{Note the recurring formula.} This is the \emph{same} $\tan\theta=v^{2}/rg$ as the conical
pendulum, the cone, and (in \S\ref{sec:roads}) the banked road, and for the same reason: here the
slanted force is the resultant of $N$ and $f$.
\end{derivbox}

\subsection{Worked examples: applying the central principle}

The examples below apply the central principle to four further circles, not all of them horizontal. In
each, the work is in deciding which real force, or which component of one, points towards the centre.

\begin{figure}[htbp]
\centering
\tikzfig{cd_groove}
\figcaption{fig:groove}{A block sliding in a circular groove, from above (left) and in section (right). The floor's push $N'$ only balances $mg$; the groove's \emph{side wall} is the only thing that can push horizontally, so the entire centripetal requirement is met by the side wall's normal force $N$.}
\end{figure}

\begin{exbox}{Block in a circular groove}
A \SI{100}{g} block moves at uniform speed in a horizontal circular groove of radius \SI{25}{cm},
completing one round in \SI{2.0}{s}. Find the normal contact force from the side wall.

\stepn{1} \textbf{Speed from the period.} One full circumference in \SI{2.0}{s}:
\[
v=\frac{2\pi r}{T_{\text{per}}}=\frac{2\pi(0.25)}{2.0}=\SI{0.785}{m/s}.
\]

\stepn{2} \textbf{Centripetal acceleration.}
\[
a=\frac{v^{2}}{r}=\frac{(0.785)^{2}}{0.25}=\SI{2.47}{m/s^2},\qquad\text{towards the centre.}
\]

\stepn{3} \textbf{Identify the force that supplies it.} The floor's normal force is vertical and
gravity is vertical; the \emph{only} horizontal force available is the push of the groove's side wall
(Figure~\ref{fig:groove}, right).
So that force alone must do the whole job:
\[
N=ma=(0.100)(2.47)=\boxed{\SI{0.247}{N}}
\]
\ans{$\approx\SI{0.25}{N}$}

\smallskip
\textbf{Note the reasoning order:} first the acceleration the motion \emph{demands}, then the real
force that pays for it. That order works every time.
\end{exbox}

\begin{figure}[htbp]
\centering
\tikzfig{cd_threeparticles}
\figcaption{fig:threepart}{Three equal masses at the corners of an equilateral triangle of side $a$. On each, the two pulls $F$ add to $\sqrt3\,F$ directed at the centroid, which is exactly where the centripetal force must point for a circular orbit of radius $r=a/\sqrt3$. (Forces are shown on one mass only.)}
\end{figure}

\begin{exbox}{Three mutually orbiting masses}
Three particles, each of mass $m$, sit at the corners of an equilateral triangle of side $a$, held
together only by their mutual gravitation. Find the speed each must have to keep the triangle rigid
while rotating, and the period.

\stepn{1} \textbf{Find the radius of the circle.} Each particle orbits the centroid. For an equilateral
triangle the centroid is a distance
\[
r=\frac{a}{2}\sec30^\circ=\frac{a}{\sqrt3}
\]
from each vertex.

\stepn{2} \textbf{Find the net force on one particle.} It is pulled by the other two, each with
$F=\dfrac{Gm^{2}}{a^{2}}$, where $G$ is the gravitational constant (not the point $G$ of the cyclist; Newton's law of gravitation, developed in \chref{ch:gravitation}; only this
inverse-square form is needed here), and the angle between those two pulls is $60^\circ$. Their resultant is
\[
F_{\text{net}}=2F\cos30^\circ=\sqrt3\,\frac{Gm^{2}}{a^{2}},
\]
directed at the centroid (Figure~\ref{fig:threepart}): exactly where a centripetal force must point. \emph{This is why the
configuration works at all.}

\stepn{3} \textbf{Set it equal to the centripetal requirement.}
\[
\sqrt3\,\frac{Gm^{2}}{a^{2}}=\frac{mv^{2}}{r}=\frac{mv^{2}\sqrt3}{a}
\qquad\Longrightarrow\qquad
\boxed{\;v=\sqrt{\frac{Gm}{a}}\;}
\]

\stepn{4} \textbf{Period.}
\[
T_{\text{per}}=\frac{2\pi r}{v}=\frac{2\pi a/\sqrt3}{\sqrt{Gm/a}}
=\boxed{\;2\pi\sqrt{\frac{a^{3}}{3Gm}}\;}
\]
\ans{$v=\sqrt{Gm/a}$;\; $T_{\text{per}}=2\pi\sqrt{a^{3}/3Gm}$}
\end{exbox}

\begin{figure}[htbp]
\centering
\tikzfig{cd_offsetarm}
\figcaption{fig:offsetarm}{When the cord hangs from an arm of length $a$ rather than from the axis
itself, the radius of the circle is $r=a+\ell\sin\theta$. Using $\ell\sin\theta$ alone (the
conical-pendulum reflex) gives the wrong answer.}
\end{figure}

\begin{exbox}{A conical pendulum on an offset arm}
The cord of the rotating apparatus in Figure~\ref{fig:offsetarm} makes $45^\circ$ with the vertical. Given
$a=\SI{20}{cm}$, $\ell=\sqrt2$\,\si{m} and $m=\SI{5.0}{kg}$, find (a) the rotation rate in rpm and (b) the
tension in the cord. Take $g=\SI{9.8}{m/s^2}$.

\stepn{1} \textbf{Get the radius right: the whole point of this example.} The bob is offset from the
axis both by the arm and by the cord:
\[
r=a+\ell\sin45^\circ=0.20+\sqrt2\left(\tfrac{1}{\sqrt2}\right)=0.20+1.00=\SI{1.20}{m}.
\]

\stepn{2} \textbf{The two equations}, exactly as for the conical pendulum:
\[
T\cos\theta=mg, \qquad T\sin\theta=mr\omega^{2}.
\]

\stepn{3} \textbf{Divide to eliminate $T$:}
\[
\tan\theta=\frac{r\omega^{2}}{g}
\qquad\Longrightarrow\qquad
\omega=\sqrt{\frac{g\tan45^\circ}{r}}=\sqrt{\frac{9.8}{1.20}}=\SI{2.86}{rad/s}.
\]

\stepn{4} \textbf{Convert to rpm.} The rotation rate $n$ in revolutions per minute is
\[
n=\frac{\omega}{2\pi}\times60=\frac{2.86\times60}{6.283}=\boxed{\SI{27.3}{rpm}}
\]

\stepn{5} \textbf{Tension}, from the vertical equation:
\[
T=\frac{mg}{\cos45^\circ}=\sqrt2\,mg=\sqrt2(5.0)(9.8)=\boxed{\SI{69.3}{N}}
\]
\ans{\SI{27.3}{rpm}; \SI{69.3}{N}}
\end{exbox}

\begin{trapbox}{the radius is that of the \emph{actual} circle}
In the plain conical pendulum $r=\ell\sin\theta$; here, with an offset arm, $r=a+\ell\sin\theta$. The formula
$T_{\text{per}}=2\pi\sqrt{\ell\cos\theta/g}$ was derived \emph{assuming} $r=\ell\sin\theta$ and is therefore
\textbf{invalid} here. When the geometry changes, go back to the two force equations rather than reaching
for the packaged result.
\end{trapbox}

\begin{exbox}{Power delivered in accelerated circular motion \textnormal{(a classic examination problem)}}
A particle of mass $m$ moves on a circle of radius $r$ with centripetal acceleration
$a_r=k^{2}rt^{2}$, where $k$ is constant. Find the power delivered to it.

\stepn{1} \textbf{Extract the speed.} By definition $a_r=v^{2}/r$, so
\[
\frac{v^{2}}{r}=k^{2}rt^{2}
\qquad\Longrightarrow\qquad
v^{2}=k^{2}r^{2}t^{2}
\qquad\Longrightarrow\qquad
v=krt.
\]
The speed grows linearly with time, so the motion is not uniform and a tangential force acts.

\stepn{2} \textbf{Power from $P=\vF\cdot\vv$.} Power (\si{W}, a scalar) is the rate at which a force does
work, $P=\vF\cdot\vv$ (quoted in \chref{ch:toolkit}, developed in \chref{ch:energy}). The net force has a
centripetal part $\vF_c$ and a tangential part $\vF_t$ of magnitude $F_t=ma_t=m\dfrac{dv}{dt}=mkr$. But
$\vF_c\perp\vv$, so it contributes \emph{nothing}:
\[
P=(\vF_c+\vF_t)\cdot\vv=F_tv=(mkr)(krt)=\boxed{\;mk^{2}r^{2}t\;}
\]

\stepn{3} \textbf{Cross-check by energy.} The power delivered equals the rate of growth of the kinetic
energy $K=\tfrac{1}{2}mv^{2}$ (the work--energy theorem of \chref{ch:energy}):
\[
P=\frac{dK}{dt}=\frac{d}{dt}\left(\tfrac{1}{2}mk^{2}r^{2}t^{2}\right)=mk^{2}r^{2}t \quad\checkmark
\]
\ans{$P=mk^{2}r^{2}t$}

\smallskip
\textbf{The lesson worth keeping:} the centripetal force \emph{never} delivers power, however large it
is, because it is always perpendicular to the velocity. Only the tangential force does work. This fact
returns in the vertical circle.
\end{exbox}

\section{Vehicles on Curved Roads}\label{sec:roads}
The cyclist showed friction on a level road doing the turning. A road can also be shaped to help: it
can supply the inward force through friction, or through the normal force on a banked surface. The three cases below take them separately and then together, and Figure~\ref{fig:roadcases} sets them side by side.

\begin{remarkbox}{A note on $\mu$ in this chapter}
Every friction limit in this chapter is a body on the \emph{verge} of slipping, so $\mu$ always means the
static coefficient $\mu_s$. A tyre that grips rolls, and its contact patch does not slide on the road;
that is why static friction is right for a car on a bend.
\end{remarkbox}

\begin{figure}[htbp]
\centering
\tikzfig{cd_roadcases}
\figcaption{fig:roadcases}{The three strategies side by side, seen from behind with the centre of the bend to the left. (a) Friction alone works but fails on a wet road; (b) banking alone works at exactly one speed, with the horizontal part of $N$ doing the turning; (c) real roads combine the two and so permit a band of safe speeds. In (c) friction is drawn for a speed above the design speed; below it, friction points up the slope.}
\end{figure}

\subsection{By friction alone (a flat road)}
On a level road the normal force is vertical and contributes nothing sideways, so friction between tyre
and road must do the entire job of turning the car. We ask how fast the car may go before friction runs
out.

\begin{derivbox}{The skid speed}
A car of mass $m$ rounds a level bend of radius $r$. Friction must do all the turning:
\[
f=\frac{mv^{2}}{r},
\qquad\text{with the ceiling}\qquad
f\le\mu N=\mu mg.
\]
Combining,
\[
\frac{mv^{2}}{r}\le\mu mg
\qquad\Longrightarrow\qquad
\boxed{\;v\le\sqrt{\mu rg}\;}
\]
Beyond this the demand exceeds what friction can supply and the car skids \emph{outward}. Mass cancels
again.

\smallskip
\textbf{Why this is unsatisfactory in practice:} it stakes everything on $\mu$, which falls sharply in
rain, ice or oil; that is why engineers bank roads.
\end{derivbox}

\subsection{By banking alone (a smooth banked road)}
Now raise the outer edge of the road, as is done on motorway curves, velodromes and racetracks, and
suppose for the moment that the surface is perfectly smooth. It is not obvious that the car can turn at
all, since friction was the only horizontal force in the previous case. The resolution is that tilting
the road also tilts the \emph{normal force}, giving it a horizontal component.

\begin{derivbox}{The design speed}
Bank the road at $\theta$ to the horizontal. The normal force, perpendicular to the surface, then makes
$\theta$ with the vertical, and its horizontal part points towards the centre of the bend.

\stepn{1} \textbf{Vertical equilibrium:} $N\cos\theta=mg$.

\stepn{2} \textbf{Radial:} $N\sin\theta=\dfrac{mv^{2}}{r}$.

\stepn{3} \textbf{Divide:}
\[
\boxed{\;\tan\theta=\frac{v^{2}}{rg}\qquad\text{or}\qquad v=\sqrt{rg\tan\theta}\;}
\]


\smallskip
\textbf{Read what this means.} For a given bank angle there is exactly \emph{one} speed, the
\textbf{design speed}, at which no friction is required. Below it, $N\sin\theta$ is more
inward force than the turn needs, so the vehicle tends to slide \emph{down} the bank; above it, too
little, so it tends to drift \emph{up} and out. The angle is independent of mass, which is why one
banking suits cars, lorries and motorcycles alike.
\end{derivbox}

\subsection{By friction and banking together (the real case)}
Real roads are banked \emph{and} rough, and vehicles travel them at many speeds. Friction is now the
awkward force, because both its size and its direction depend on the speed. Rather than a single safe speed there
is now a \emph{band} of them, and we want its two edges.

\begin{motivbox}{How friction knows which way to point}
Friction's direction is decided by the vehicle's speed relative to the design speed
$v_0=\sqrt{rg\tan\theta}$:
\begin{itemize}
\item $v<v_0$: the vehicle tends to slide \emph{down} the slope, so friction acts \textbf{up} the slope.
\item $v=v_0$: no tendency either way, so $f=0$: the banking alone is doing the whole job.
\item $v>v_0$: the vehicle tends to slide \emph{up} and out, so friction acts \textbf{down} the slope.
\end{itemize}
As in \chref{ch:friction}, friction self-adjusts up to its ceiling $\mu N$; the two limits below are
where it reaches that ceiling in each direction.
\end{motivbox}

\begin{derivbox}{Deriving the limiting speeds}
\stepn{1} \textbf{Fastest case.} The vehicle tends to slide up and out, so friction acts \emph{down} the
slope at its limit $\mu N$. Its components are $\mu N\cos\theta$ inward and $\mu N\sin\theta$ downward.
\[
\text{Vertical: } N\cos\theta-\mu N\sin\theta=mg,
\qquad
\text{Radial: } N\sin\theta+\mu N\cos\theta=\frac{mv^{2}}{r}.
\]
\stepn{2} \textbf{Divide} to eliminate $N$ and $m$, then divide top and bottom by $\cos\theta$:
\[
\frac{v_{\max}^{2}}{rg}=\frac{\sin\theta+\mu\cos\theta}{\cos\theta-\mu\sin\theta}
=\frac{\tan\theta+\mu}{1-\mu\tan\theta}.
\]
\stepn{3} \textbf{Slowest case.} Friction now acts \emph{up} the slope, which is the same working with
$\mu$ replaced by $-\mu$:
\[
\frac{v_{\min}^{2}}{rg}=\frac{\tan\theta-\mu}{1+\mu\tan\theta}.
\]
\end{derivbox}

\begin{keybox}{The two limiting speeds}
Resolving with friction at its limit $\mu N$, acting down the slope (fastest case) or up it (slowest):
\[
\boxed{\;v_{\max}=\sqrt{rg\,\frac{\tan\theta+\mu}{1-\mu\tan\theta}}\;}
\qquad\qquad
\boxed{\;v_{\min}=\sqrt{rg\,\frac{\tan\theta-\mu}{1+\mu\tan\theta}}\;}
\]
\textbf{Checks.} Setting $\mu=0$ collapses both to the design speed $\sqrt{rg\tan\theta}$ \checkmark.
If $\mu\ge\tan\theta$ the expression under the root for $v_{\min}$ is zero or negative: there is no
lower limit, and the vehicle can stand still on the bank without sliding. And $v_{\max}$ grows without
bound as $\mu\tan\theta\to1$; for $\mu\tan\theta\ge1$ the bank is so steep and so rough that no
speed makes the vehicle slide up it (in this model).
\end{keybox}

\subsection{Worked example: which way does friction point?}

The example below works through both sides of the design speed, where friction points in opposite
directions (Figure~\ref{fig:fricdir}).

\begin{figure}[htbp]
\centering
\tikzfig{cd_frictiondir}
\figcaption{fig:fricdir}{On a banked road, friction is a self-adjusting force whose \emph{direction}
flips at the design speed. Below $v_0$ it acts up the slope to stop the vehicle sliding down; above
$v_0$ it acts down the slope to stop it sliding up and out.}
\end{figure}

\begin{exbox}{Friction on a banked turn, above and below the design speed}
A turn of radius \SI{20}{m} is banked for a \SI{200}{kg} vehicle travelling at \SI{10}{m/s}. Find the
direction and magnitude of the friction when the vehicle instead moves at (a) \SI{5}{m/s},
(b) \SI{15}{m/s}. Take $g=\SI{10}{m/s^2}$ and assume friction is sufficient.

\stepn{1} \textbf{Find the banking angle from the design speed.}
\[
\tan\theta=\frac{v_0^{2}}{rg}=\frac{10^{2}}{20\times10}=\frac12.
\]

\stepn{2} \textbf{(a) At \SI{5}{m/s}, slower than designed.} The vehicle tends to slide \emph{down}
the bank, so friction acts \textbf{up} the slope. Resolving horizontally (towards the centre) and
vertically:
\[
N\sin\theta-f\cos\theta=\frac{mv^{2}}{r},
\qquad
N\cos\theta+f\sin\theta=mg.
\]
Rather than solve these simultaneously, resolve \emph{along} the slope, where $N$ does not appear.
The inward acceleration is horizontal and the slope makes $\theta$ with the horizontal, so the
acceleration has a component $(v^{2}/r)\cos\theta$ down the slope, while the weight has $mg\sin\theta$. With
$\sin\theta=1/\sqrt5$, $\cos\theta=2/\sqrt5$, $mg=\SI{2000}{N}$ and $mv^{2}/r=\SI{250}{N}$:
\[
mg\sin\theta-f=\frac{mv^{2}}{r}\cos\theta
\;\Longrightarrow\;
f=\frac{2000-500}{\sqrt5}=\boxed{300\sqrt5\approx\SI{671}{N}\ \text{up the slope}}
\]

\stepn{3} \textbf{(b) At \SI{15}{m/s}, faster than designed.} Now the vehicle tends to slide
\emph{up} and outward, so friction reverses to act \textbf{down} the slope. Resolving along the slope
as before, now with $mv^{2}/r=\SI{2250}{N}$:
\[
f=\frac{mv^{2}}{r}\cos\theta-mg\sin\theta=\frac{4500-2000}{\sqrt5}
=\boxed{500\sqrt5\approx\SI{1118}{N}\ \text{down the slope}}
\]
\ans{(a) $300\sqrt5$\,N up; (b) $500\sqrt5$\,N down}

\smallskip
\textbf{Do not memorise the sign pattern; derive it.} Ask only: \emph{if friction vanished, which way
would the vehicle slide?} Friction then points the other way.
\end{exbox}

\subsection{The second critical speed: toppling}
Everything so far has asked when the vehicle \emph{slides}. But a tall, narrow vehicle such as a loaded
lorry or a double-decker bus can roll over while its tyres still grip, and the skidding analysis gives
no warning of this because friction never reaches its limit. There is a \emph{second} critical speed,
governed by the vehicle's shape, and the one that matters is whichever of the two is smaller.

\begin{motivbox}{Skidding is not the only way to come off}
Friction acts at ground level, but the weight acts at the centre of gravity, a height $h$ up. That offset
gives friction a \emph{torque} about the centre of gravity, which transfers load from the inner wheels
to the outer ones. When the inner wheels carry nothing at all, the vehicle tips.
\end{motivbox}

\begin{figure}[htbp]
\centering
\tikzfig{cd_topple}
\figcaption{fig:topple}{Seen from behind, with the centre of the bend to the left. Toppling. The wheels are $2a$ apart and the centre of gravity is at height $h$.
Friction $f$ at road level has a moment about $G$ that unloads the inner wheel ($N_2$) and loads the
outer one ($N_1$). When $N_2$ reaches zero, the vehicle overturns.}
\end{figure}

\begin{derivbox}{The toppling speed}
Take a vehicle of mass $m$ on a level bend of radius $r$, with $N_1$ and $N_2$ the total normal
reactions on the outer and inner wheels and $f$ the total friction (Figure~\ref{fig:topple}).

\stepn{1} \textbf{Vertical equilibrium:} $N_1+N_2=mg$.

\stepn{2} \textbf{Radial:} friction supplies the turn, $f=\dfrac{mv^{2}}{r}$.

\stepn{3} \textbf{Rotational equilibrium about $G$.} The vehicle is not rotating, so moments about the
centre of gravity must balance. (As for the cyclist, take moments about $G$ itself: the vehicle is
accelerating, and about any other point the balance would need an extra term.) With the wheels $2a$
apart and $G$ at height $h$, the outer reaction $N_1$ (lever arm $a$) turns the vehicle one way, and the
inner reaction $N_2$ (arm $a$) and friction (arm $h$) turn it the other:
\[
N_1 a = N_2 a + f h.
\]
\textbf{Read this:} friction's moment $fh$ is what forces $N_1$ to exceed $N_2$.

\stepn{4} \textbf{Solve for the inner reaction.} Substituting $N_1=mg-N_2$ and $f=mv^{2}/r$,
\[
N_2=\frac{mg}{2}-\frac{mv^{2}}{r}\cdot\frac{h}{2a}.
\]
The faster you go, the smaller $N_2$ becomes: the inner wheel is progressively unloaded.

\stepn{5} \textbf{The limit.} Setting $N_2=0$:
\[
\boxed{\;v_{\text{topple}}=\sqrt{\frac{gra}{h}}\;}
\]


\smallskip
\textbf{The model's limits.} This treats the vehicle as rigid on a flat road. A real vehicle rolls on its
suspension, which shifts $G$ outward and lowers the true toppling speed below this estimate.
\end{derivbox}

\begin{remarkbox}{Which speed matters, and what engineers do about it}
\[
v_{\text{skid}}=\sqrt{\mu rg},
\qquad\qquad
v_{\text{topple}}=\sqrt{\frac{gra}{h}}.
\]
The vehicle fails at whichever is \textbf{lower}, so safe driving means staying below both.
\begin{itemize}
\item $v_{\text{skid}}\propto\sqrt{\mu}$: better tyres and road surface raise it.
\item $v_{\text{topple}}\propto\sqrt{a/h}$: a \emph{wide} track ($a$ large) and a \emph{low} centre of
      gravity ($h$ small) raise it. This is why racing cars are wide and low, and why a lorry piled
      high with sugarcane (which raises $h$) can overturn at surprisingly low speeds.
\end{itemize}
\textbf{Designers prefer skidding to toppling}, because a skid can often be recovered and a rollover
generally cannot. So a vehicle is built, where possible, with $v_{\text{skid}}<v_{\text{topple}}$,
that is, with $\mu<a/h$.
\end{remarkbox}

\section{Motion in a Vertical Circle}\label{sec:vertcircle}

In every horizontal circle above, gravity was perpendicular to the motion and did no work, so the speed
could stay constant. Hang the circle vertically and that comfort disappears. Now gravity has a component
\emph{along} the path: it helps on the way down and fights on the way up. So the speed changes
continuously, fastest at the bottom, slowest at the top, which means
\[
\text{this is \textbf{non-uniform} circular motion: both } a_t \text{ and } a_r \text{ are non-zero.}
\]
Two consequences follow, and they dictate the whole method.

\smallskip
\textbf{First:} one equation is no longer enough. Newton's law along the radius gives the tension at a
point only if the speed there is known, and the speed varies. So \emph{energy conservation} supplies the
speed, and \emph{Newton's radial equation} converts it into a tension. \textbf{A typical
vertical-circle problem uses both laws, in that order.}

\smallskip
\textbf{Second:} a string can only \emph{pull}, never push. At the top, gravity already points towards
the centre. If gravity alone is more than enough to bend the path, the string would have to push, which
it cannot; it goes \textbf{slack}, and the particle leaves the circle.

\subsection{The two master equations}
Consider a small bob of mass $m$ tied to a light inextensible string, swung in a vertical circle and
given a speed $u$ at the lowest point. Nothing drives it after that; it coasts. While the string is taut
its length is also the radius of the circle, so in this section's theory we call it $R$ (the worked
examples use $\ell$). We want the \emph{speed} and the \emph{tension} at an arbitrary position,
specified by the angle $\theta$ turned from the bottom (Figure~\ref{fig:vcgen}); every other result in the section follows from
these two.

\begin{figure}[htbp]
\centering
\tikzfig{vc_general}
\figcaption{fig:vcgen}{The general position, at angle $\theta$ from the lowest point. The height risen, $h=R(1-\cos\theta)$, fixes the speed through energy conservation. The radial balance then fixes the tension: $T$ acts inwards, while the weight's component $mg\cos\theta$ acts outwards along the radius. Its other component, $mg\sin\theta$, lies along the tangent and slows the bob.}
\end{figure}

\begin{derivbox}{Speed and tension at any angle}
With $\theta$ measured from the lowest point, the particle has risen a height $h$ (a new use of the
letter: not the height of $G$ in \S\ref{sec:roads})
\[
h=R-R\cos\theta=R(1-\cos\theta). \tag{geometry}
\]

\stepn{1} \textbf{Energy conservation gives the speed.} The string is always perpendicular to the
motion, so the tension does \emph{no work}; only gravity does. Hence
\[
\tfrac{1}{2}mv^{2}=\tfrac{1}{2}mu^{2}-mgh
\qquad\Longrightarrow\qquad
\boxed{\;v^{2}=u^{2}-2gh\;} \tag{I}
\]
Energy conservation is developed properly in \chref{ch:energy}. The same result follows from Newton's
law alone. Along the tangent the only force is the component $-mg\sin\theta$ of the weight, so
$m\dfrac{dv}{dt}=-mg\sin\theta$. With arc length $s=R\theta$, the chain rule gives
$\dfrac{dv}{dt}=\dfrac{dv}{ds}\dfrac{ds}{dt}=\dfrac{v}{R}\dfrac{dv}{d\theta}$, so
$v\,dv=-gR\sin\theta\,d\theta$. Integrating from $(0,u)$ to $(\theta,v)$ gives
$v^{2}=u^{2}-2gR(1-\cos\theta)$, which is (I).

\stepn{2} \textbf{Newton's law along the radius gives the tension.} Take the positive direction
\textbf{inward}, towards the centre. The tension acts inward along the string, so it is $+T$. The
weight's component along the inward radius is $-mg\cos\theta$: at the bottom ($\theta=0$) the weight
points directly away from the centre, giving $-mg$; at the top ($\theta=180^\circ$) it points directly
towards the centre, and $\cos180^\circ=-1$ turns the term into $+mg$. One expression therefore covers
the whole circle:
\[
T-mg\cos\theta=\frac{mv^{2}}{R}
\qquad\Longrightarrow\qquad
T=mg\cos\theta+\frac{mv^{2}}{R}. \tag{II}
\]

\stepn{3} \textbf{Combine.} Substituting (I) with $h=R(1-\cos\theta)$ into (II),
$T=mg\cos\theta+\dfrac{mu^{2}}{R}-2mg(1-\cos\theta)$, which collects to
\[
\boxed{\;T=\frac{mu^{2}}{R}+mg\,(3\cos\theta-2)\;}
\qquad\text{equivalently}\qquad
T=mg+\frac{mu^{2}}{R}-\frac{3mgh}{R}.
\]

\stepn{4} \textbf{Read the result before using it.} The tension \emph{falls} steadily as the particle
rises ($\cos\theta$ decreases, $h$ increases). So $T$ is
\textbf{greatest at the bottom} ($\theta=0$) and \textbf{least at the top} ($\theta=180^\circ$):
\[
T_{\text{bottom}}=\frac{mu^{2}}{R}+mg,
\qquad\qquad
T_{\text{top}}=\frac{mu^{2}}{R}-5mg.
\]
\end{derivbox}

\why{\textbf{Why the tension is largest at the bottom.} There the string must support the weight
\emph{and} supply the inward force; at the top gravity \emph{helps}. The particle is also faster at the
bottom. So a string whirled in a vertical circle, if it snaps, snaps at the bottom.}

\begin{keybox}{A result worth remembering: $T_b-T_t=6mg$}
Subtracting the two tensions of Step~4:
\[
T_{\text{bottom}}-T_{\text{top}}=\left(\frac{mu^{2}}{R}+mg\right)-\left(\frac{mu^{2}}{R}-5mg\right)
=\boxed{6mg}
\]
This holds for \emph{any} launch speed that completes the circle. Of the $6mg$, $2mg$ arises because the
weight switches from opposing the tension to assisting it, and $4mg$ from the speed difference, since energy
conservation gives $v_{\text{bottom}}^{2}-v_{\text{top}}^{2}=4gR$ across the height $2R$.
\end{keybox}

\subsection{The critical condition for completing the circle}
Give the bob too little speed at the bottom and it will not get round; give it enough and it will.
We now find the critical launch speed between the two.

\begin{derivbox}{$u_{\min}=\sqrt{5gR}$}
\stepn{1} \textbf{Where is the danger point?} The tension falls steadily as the bob rises, so it is
smallest at the very top. If the string goes slack anywhere, it goes slack there first, and the top is
the only point we need to test.

\stepn{2} \textbf{Impose the physical constraint.} A string cannot push, so $T\ge0$. At the top:
\[
T_{\text{top}}=\frac{mu^{2}}{R}-5mg\ \ge\ 0
\qquad\Longrightarrow\qquad
\boxed{\;u\ge\sqrt{5gR}\;}
\]

\stepn{3} \textbf{What the critical case looks like.} At exactly $u=\sqrt{5gR}$:
\[
T_{\text{top}}=0,\qquad
v_{\text{top}}=\sqrt{u^{2}-4gR}=\sqrt{gR},\qquad
T_{\text{bottom}}=6mg.
\]
At that instant the string is limp and \textbf{gravity alone} supplies the entire centripetal
requirement: $mg=mv_{\text{top}}^{2}/R$ gives $v_{\text{top}}=\sqrt{gR}$. Zero tension does
\emph{not} mean the particle drops straight down: it is moving horizontally at $\sqrt{gR}$, and a
projectile with that speed has its acceleration $g$ perpendicular to its velocity, so its path has radius
of curvature $v^{2}/g=R$ (\S\ref{sec:curvature}): exactly the circle's. Momentarily the particle is in
free fall, yet still on the circle.
\end{derivbox}

\begin{exbox}{A comfortable margin}
A particle is given $u=\sqrt{8gR}$ at the bottom. Find the speed and tension at the top and bottom.

\stepn{1} \textbf{Top speed:} $v_t^{2}=u^{2}-4gR=8gR-4gR=4gR$, so $v_t=2\sqrt{gR}$.

\stepn{2} \textbf{Top tension:} $T_t=\dfrac{mu^{2}}{R}-5mg=8mg-5mg=\mathbf{3mg}$, comfortably
positive, so the string stays taut.

\stepn{3} \textbf{Bottom tension:} $T_b=\dfrac{mu^{2}}{R}+mg=8mg+mg=\mathbf{9mg}$.

\stepn{4} \textbf{Check:} $T_b-T_t=9mg-3mg=6mg$ \checkmark, as it must be.
\ans{$v_t=2\sqrt{gR}$; $T_t=3mg$; $T_b=9mg$}
\end{exbox}

\subsection{What happens when the particle cannot complete the circle}\label{sub:vcfail}
Suppose $u<\sqrt{5gR}$, so the bob does not get round. Many students assume it simply swings back
like a pendulum, but that is only one of two possible fates; Figure~\ref{fig:vertcases} shows both beside the complete circle.

\begin{motivbox}{Two different ways to fail, and the quantity that decides}
Which fate occurs depends on \emph{what runs out first} as the particle climbs:
\begin{itemize}
\item its \textbf{speed} runs out first ($v\to0$ while the string is still taut): the particle
      stops, swings back, and \textbf{oscillates} like a pendulum;
\item its \textbf{tension} runs out first ($T\to0$ while it is still moving): the string goes slack
      and the particle \textbf{leaves the circle}, becoming a projectile.
\end{itemize}
So compute the height at which each would happen, and see which comes sooner.
\end{motivbox}

\begin{derivbox}{The two critical heights}
\stepn{1} \textbf{Height $h_1$ where the speed vanishes.} Put $v=0$ in (I):
\[
0=u^{2}-2gh_1
\qquad\Longrightarrow\qquad
\boxed{\;h_1=\frac{u^{2}}{2g}\;}
\]

\stepn{2} \textbf{Height $h_2$ where the tension vanishes.} Put $T=0$ in the combined formula
$T=mg+\dfrac{mu^{2}}{R}-\dfrac{3mgh}{R}$:
\[
0=mg+\frac{mu^{2}}{R}-\frac{3mgh_2}{R}
\qquad\Longrightarrow\qquad
\boxed{\;h_2=\frac{u^{2}+gR}{3g}\;}
\]

\stepn{3} \textbf{Compare them.} The particle meets whichever height is \emph{lower}.
\[
h_2<h_1
\iff
\frac{u^{2}+gR}{3g}<\frac{u^{2}}{2g}
\iff
2u^{2}+2gR<3u^{2}
\iff
\boxed{u^{2}>2gR}
\]
So:
\begin{itemize}
\item If $u>\sqrt{2gR}$: tension dies first $\Rightarrow$ the particle \textbf{leaves the circle}.
\item If $u\le\sqrt{2gR}$: speed dies first $\Rightarrow$ the particle \textbf{oscillates}.
\end{itemize}
\end{derivbox}

\begin{figure}[htbp]
\centering
\tikzfig{cd_vertcases}
\figcaption{fig:vertcases}{The three fates, decided entirely by the launch speed $u$. Completing the
circle needs $u\ge\sqrt{5gR}$; below $\sqrt{2gR}$ the particle merely swings; in between it flies off
the circle partway up the far side and continues as a projectile.}
\end{figure}

\begin{keybox}{The three regimes in full}
\begin{enumerate}
\item \textbf{$u\ge\sqrt{5gR}$, completes the circle.} The string stays taut all the way round.
\item \textbf{$\sqrt{2gR}<u<\sqrt{5gR}$, leaves the circle.} The string goes slack at the height
      $h_2=\dfrac{u^{2}+gR}{3g}$, which lies in the \emph{upper half}
      ($90^\circ<\theta<180^\circ$), since $h_2>R$ exactly when $u^{2}>2gR$. From that moment there is no string force at all, so the particle
      becomes an ordinary \textbf{projectile}, launched along the tangent with the speed it had, and
      follows a parabola inside the circle until the string snaps taut again.
\item \textbf{$0<u\le\sqrt{2gR}$, oscillates.} The particle rises to $h_1=\dfrac{u^{2}}{2g}$, which
      lies in the \emph{lower half} ($0^\circ<\theta\le90^\circ$, since $h_1\le R$), momentarily stops with the string
      still taut, and swings back. This is exactly a pendulum.
\end{enumerate}
\textbf{The boundary case} $u=\sqrt{2gR}$ gives $h_1=h_2=R$: speed and tension vanish together, exactly
at the horizontal level of the centre.
\end{keybox}

\begin{remarkbox}{A rigid rod changes the answer, and here is why}
Replace the string with a light \textbf{rod}. A rod can \emph{push} as well as pull, so the constraint
$T\ge0$ disappears: at the top the rod can support the particle from below. The only remaining
requirement is that the particle actually \emph{pass over} the top, that is, $v_{\text{top}}>0$:
\[
u^{2}-4gR>0
\qquad\Longrightarrow\qquad
\boxed{\;u>2\sqrt{gR}\;}
\]
At exactly $2\sqrt{gR}$ the particle only creeps ever closer to the top, so this is a limit, not an
attainable value. Compare this with
$\sqrt{5gR}$ for a string. The same relaxation applies to a bead threaded on a circular wire or a
particle inside a smooth tube, anything that can push inward \emph{or} outward. By contrast, for a
particle on the \emph{inside} of a smooth circular track or spherical shell, the normal force $N$
plays exactly the role of the string's tension (it can only push inward), so the string conditions apply
unchanged.
\end{remarkbox}

\subsection{A ball sliding off a smooth sphere}\label{sub:sphere}
Every vertical circle so far had the body on the \emph{inside}, with the support pulling or pushing it
towards the centre. Put the body on the outside and the same two laws give a different story. A ball balanced at the top of a large smooth dome is given the gentlest of nudges and slides down. It
does not hug the surface all the way to the bottom; at some point it flies off. The question is
\emph{where}, and the answer is a pure number, independent of the dome's size, the ball's mass and even
the strength of gravity.

\begin{motivbox}{The mirror image of the loop}
On the outside of the curve the normal force points \emph{away} from the centre, while gravity points
inward. Now it is gravity that supplies the turning and the
normal force that opposes it. As the ball speeds up, more and more of gravity is consumed by the
curving, leaving less for pressing on the surface. When nothing is left, $N=0$ and contact is lost.
A car cresting a hump-backed bridge of radius $R$ is the everyday case: at the top $N=mg-mv^{2}/R$, so
above $\sqrt{gR}$ the wheels leave the road.
\end{motivbox}

\begin{figure}[htbp]
\centering
\tikzfig{cd_sphere}
\figcaption{fig:sphere}{A ball released from rest at the top of a smooth sphere. Here the normal force $N$ points outwards, away from the centre, so it opposes the inward pull of gravity. As the ball speeds up, $N$ falls; it reaches zero at $\cos\theta=\tfrac{2}{3}$, after a drop of $R/3$. From there the ball flies off as a projectile.}
\end{figure}

\begin{derivbox}{Where the ball leaves, $\cos\theta=\tfrac{2}{3}$}
A ball starts from rest at the top of a smooth sphere of radius $R$ and slides down. Let $\theta$ be the
angle turned from the top (Figure~\ref{fig:sphere}).

\stepn{1} \textbf{Radial equation.} Towards the centre, the inward component of gravity minus the
outward normal force gives
\[
mg\cos\theta-N=\frac{mv^{2}}{R}. \tag{i}
\]

\stepn{2} \textbf{Energy conservation.} It has descended a height $h=R(1-\cos\theta)$:
\[
v^{2}=2gR(1-\cos\theta). \tag{ii}
\]

\stepn{3} \textbf{Impose the leaving condition} $N=0$ and substitute (ii) into (i):
\[
mg\cos\theta=\frac{m}{R}\cdot2gR(1-\cos\theta)
\qquad\Longrightarrow\qquad
\cos\theta=2-2\cos\theta
\]
\[
\Longrightarrow\qquad
3\cos\theta=2
\qquad\Longrightarrow\qquad
\boxed{\;\cos\theta=\frac{2}{3},\qquad \theta=48.2^\circ\;}
\]
The answer is independent of $R$, of $m$ and of $g$.

\stepn{4} \textbf{If the ball is given an initial speed $u$ at the top}, (ii) becomes
$v^{2}=u^{2}+2gR(1-\cos\theta)$ and the same steps give
\[
\cos\theta=\frac{2}{3}+\frac{u^{2}}{3gR}.
\]
With $u=\tfrac{1}{2}\sqrt{gR}$ this is $\cos\theta=\tfrac{3}{4}$, that is, $\theta=41.4^\circ$, so it leaves
\emph{earlier}, as a faster start should imply. And if $u\ge\sqrt{gR}$ the formula demands
$\cos\theta\ge1$, meaning the ball leaves \emph{immediately} at the top and never touches the sphere at
all.

\smallskip
\textbf{The model's limits.} This assumes a smooth surface, so the ball slides without spinning. A
uniform solid ball rolling without slipping on a rough dome keeps some of its energy in rotation, moves
more slowly at each angle, and leaves later, at $\cos\theta=\tfrac{10}{17}$ ($\theta\approx54^\circ$;
rolling is treated in \chref{ch:rolling}).
\end{derivbox}

\begin{trapbox}{inside a loop versus outside a sphere}
These two problems look alike and have opposite signs. \textbf{Inside} a loop (or on a string) the
support force points \emph{towards} the centre and \emph{adds} to the centripetal total. \textbf{Outside}
a sphere it points \emph{away} and \emph{subtracts}. Draw the normal force before writing anything: it
always points from the surface \emph{into} the body.
\end{trapbox}

\subsection{Worked examples: vertical circles}

Five examples follow. The first needs only the radial equation at one point, because the speed there is
given; the others combine the radial equation at a point with energy conservation between points.

\begin{figure}[htbp]
\centering
\tikzfig{cd_fighter}
\figcaption{fig:fighter}{At the lowest point of a dive, the centre of the circular path is directly \emph{above} the aircraft. So ``towards the centre'' means \emph{upwards}, and the air's force $F$ must exceed the weight by exactly $mv^{2}/r$. The two arrows are drawn to the worked example's ratio, $F\approx4.2\,mg$.}
\end{figure}

\begin{exbox}{A fighter plane pulling out of a dive}
A \SI{16000}{kg} fighter pulls out of a dive at \SI{900}{km/h} on a vertical circle of radius
\SI{2000}{m}. Find the vertical force the air exerts on it at the lowest point. Take
$g=\SI{9.8}{m/s^2}$.

\stepn{1} \textbf{Convert the speed.}
$v=\SI{900}{km/h}=\dfrac{900\times1000}{3600}=\SI{250}{m/s}$.

\stepn{2} \textbf{Which way is ``towards the centre''?} At the \emph{lowest} point of the loop the centre
is directly above, so the centripetal direction is \textbf{up} (Figure~\ref{fig:fighter}).

\stepn{3} \textbf{Newton's law along that direction.} Two forces act: the air's force $F$ (up) and the
weight $mg$ (down):
\[
F-mg=\frac{mv^{2}}{r}
\qquad\Longrightarrow\qquad
F=m\left(g+\frac{v^{2}}{r}\right).
\]

\stepn{4} \textbf{Substitute.}
\[
F=16000\left(9.8+\frac{250^{2}}{2000}\right)=16000\,(9.8+31.25)=\boxed{\SI{6.57e5}{N}}
\]
\ans{\SI{6.57e5}{N}, upward}

\smallskip
\textbf{The physical point:} this is about $4.2$ times the plane's weight, and the $v^{2}/r$ term, not
gravity, dominates it. The seat must push the pilot up with the same factor, ``4.2\,$g$'', which is why
pulling out of a dive too sharply can cause the pilot to black out.
\end{exbox}

\begin{exbox}{Change in velocity, not change in speed \textnormal{(a classic examination problem)}}
A stone on a string of length $\ell$ is whirled in a vertical circle and has speed $u$ at the lowest point.
Find the \emph{magnitude of the change in its velocity} by the time the string is horizontal. (Assume
$u^{2}\ge2g\ell$, so the stone reaches the horizontal with the string still taut.)

\stepn{1} \textbf{Speed when horizontal.} It has risen $h=\ell$, so by energy conservation
\[
v^{2}=u^{2}-2g\ell.
\]

\stepn{2} \textbf{The directions: this is the whole question.} At the lowest point the velocity is
\emph{horizontal}, of magnitude $u$. When the string is horizontal the velocity is \emph{vertical}, of
magnitude $v$. The two are therefore at $90^\circ$ to each other.

\stepn{3} \textbf{Subtract as vectors.} Since the final and initial velocities are perpendicular, $\vv_f\perp\vv_i$,
\[
|\Delta\vv|=|\vv_f-\vv_i|=\sqrt{v^{2}+u^{2}}=\sqrt{(u^{2}-2g\ell)+u^{2}}
\]
\[
\boxed{\;|\Delta\vv|=\sqrt{2(u^{2}-g\ell)}\;}
\]
\ans{$\sqrt{2(u^{2}-g\ell)}$}

\smallskip
\textbf{The trap this question is built on} is answering $u-v$. That is the change in \emph{speed}; the
question asked for the change in \emph{velocity}, and the direction turned through $90^\circ$. Whenever
a question says ``change in velocity'', draw the two arrows and subtract them as vectors.
\end{exbox}

\begin{figure}[htbp]
\centering
\tikzfig{vc_slack}
\figcaption{fig:vcslack}{Example~\ref{ex:slack}. The string goes slack at the point where the tension
vanishes, necessarily in the upper half of the circle, after which the stone is an ordinary
projectile and rises a little further under gravity alone.}
\end{figure}

\begin{exbox}{Where the string slackens, and how much higher it goes}\label{ex:slack}
A particle hangs from a fixed point by a string of length \SI{5}{m}. It is projected from the lowest
point so that the string goes slack after the particle has risen \SI{8}{m}. Find the speed just before
the string slackens, and how much higher the particle then rises. Take $g=\SI{9.8}{m/s^2}$.

\stepn{1} \textbf{Locate the point geometrically} (Figure~\ref{fig:vcslack}). With $h=\ell(1-\cos\theta)$ and $h=8$, $\ell=5$:
\[
\cos\theta=\frac{\ell-h}{\ell}=\frac{5-8}{5}=-0.6,
\]
so $\theta\approx126.9^\circ$, in the \emph{upper} half, exactly as the theory of \S\ref{sub:vcfail} predicts
for a string that slackens.

\stepn{2} \textbf{Use $T=0$ in the radial equation.} At that point gravity alone supplies the
centripetal force; its component along the inward radius is $mg|\cos\theta|$, so
\[
mg|\cos\theta|=\frac{mv^{2}}{\ell}
\qquad\Longrightarrow\qquad
v^{2}=g\ell|\cos\theta|=9.8(5)(0.6)=29.4
\]
\[
v=\boxed{\SI{5.42}{m/s}}
\]

\stepn{3} \textbf{After the string slackens it is a projectile.} Its velocity is along the tangent,
perpendicular to the string; the string makes $180^\circ-\theta$ with the vertical, so the velocity makes
that angle with the horizontal, and its sine is $\sin\theta=0.8$. Only the vertical component is lost to
gravity, so the further rise is
\[
h'=\frac{(v\sin\theta)^{2}}{2g}=\frac{(5.42)^{2}(0.8)^{2}}{2(9.8)}
=\frac{29.4\times0.64}{19.6}=\boxed{\SI{0.96}{m}}
\]
\ans{$v=\SI{5.42}{m/s}$; rises a further \SI{0.96}{m}}
\end{exbox}

\begin{exbox}{Finding the angle at which $T=mg$}
A heavy particle hangs from a fixed point by an inextensible string of length $\ell$ and is projected
horizontally with speed $\sqrt{g\ell}$. Find its speed and the string's inclination to the vertical at the
moment the tension equals the particle's weight.

\stepn{1} \textbf{Energy conservation} with $h=\ell(1-\cos\theta)$ and $u^{2}=g\ell$:
\[
v^{2}=u^{2}-2g\ell(1-\cos\theta)=g\ell-2g\ell(1-\cos\theta). \tag{i}
\]

\stepn{2} \textbf{Radial equation with the stated condition $T=mg$:}
\[
T-mg\cos\theta=\frac{mv^{2}}{\ell}
\quad\Longrightarrow\quad
mg-mg\cos\theta=\frac{mv^{2}}{\ell}
\quad\Longrightarrow\quad
v^{2}=g\ell(1-\cos\theta). \tag{ii}
\]

\stepn{3} \textbf{Equate (i) and (ii).}
\[
g\ell-2g\ell(1-\cos\theta)=g\ell(1-\cos\theta)
\]
\[
1-2+2\cos\theta=1-\cos\theta
\qquad\Longrightarrow\qquad
3\cos\theta=2
\qquad\Longrightarrow\qquad
\boxed{\cos\theta=\tfrac{2}{3},\quad\theta=48.2^\circ}
\]

\stepn{4} \textbf{Speed, from (ii):}
\[
v^{2}=g\ell\left(1-\tfrac{2}{3}\right)=\frac{g\ell}{3}
\qquad\Longrightarrow\qquad
\boxed{v=\sqrt{\frac{g\ell}{3}}}
\]
\textbf{Check that the particle gets there.} With $u^{2}=g\ell$ it would stop at
$h_1=u^{2}/2g=\ell/2$, that is at $\cos\theta=\tfrac12$ ($60^\circ$); since $48.2^\circ<60^\circ$ the
string is still taut and the particle still moving at the required point.
\ans{$v=\sqrt{g\ell/3}$ at $\theta=\cos^{-1}(2/3)=48.2^\circ$}

\end{exbox}

\begin{exbox}{A tube can push, so the rules relax}
A small ball is inside a smooth circular \emph{tube} of radius $R$ held in a vertical plane, and is
pushed along the tube with speed $u$ from a point at height $h$ above the lowest point. What speed is
needed for it to pass over the top? The tube's bore is negligible compared with $R$.

\stepn{1} \textbf{Recognise what kind of constraint a tube is.} A tube has walls on \emph{both} sides,
so it can push the ball inward \emph{or} outward. The $T\ge0$ condition that governs a string therefore
does not apply: the only requirement is that the ball physically \emph{passes over} it, that is,
$v_{\text{top}}>0$.

\stepn{2} \textbf{Energy conservation} from height $h$ to the top of the tube at height $2R$:
\[
v_{\text{top}}^{2}=u^{2}-2g(2R-h).
\]

\stepn{3} \textbf{Impose the (relaxed) condition} $v_{\text{top}}>0$:
\[
u^{2}-2g(2R-h)>0
\qquad\Longrightarrow\qquad
\boxed{\;u>\sqrt{2g(2R-h)}\;}
\]
As with the rod, the boundary value is a limit the ball approaches but never attains.
\ans{$u>\sqrt{2g(2R-h)}$}

\smallskip
\textbf{Compare with a string.} For a string, starting at a height $h$ in the lower half, the
requirement would be $u^{2}\ge5gR-2gh$, larger than the tube's $4gR-2gh$ by $gR$. Always check first
\emph{what the constraint can do}: a string pulls only; a rod, a tube or a groove can push as well.
\end{exbox}

\section{Non-Inertial Frames and Pseudo-Forces}\label{sec:noninertial}

Every problem so far was solved from the ground. Yet a passenger in a turning car, or a rider in a
rotor, naturally describes the motion from where they sit, and feels pushed outward. Newton's laws
hold in some frames and not in others. This section shows how to work honestly in a frame where they
fail, first an accelerating one (a lift, a train), then a rotating one: the cost is one fictitious
force, called a pseudo-force.

\begin{defbox}{Inertial and non-inertial frames}
An \emph{inertial} frame is one in which the first law holds (\chref{ch:newton}). A frame that
accelerates or rotates relative to an inertial frame is \emph{non-inertial}: in it a free body appears
to accelerate with no real force acting, so Newton's laws fail unless they are corrected.
No real frame is exactly inertial: the ground shares the Earth's spin, whose centripetal acceleration
at the equator is about \SI{0.034}{m/s^2}, a third of a percent of $g$. That is negligible in this
book's problems, which is why the ground is treated as inertial.
\end{defbox}

\begin{derivbox}{The repair: pseudo-forces}
Suppose your frame (a train, a lift) accelerates at $\va_0$ relative to the ground \emph{without
rotating}, so that every point of it shares that acceleration. Watch a particle from
inside that frame, where it appears to accelerate at $\va'$.

\stepn{1} \textbf{Relate the two accelerations.} The particle's \emph{true} acceleration (as seen from the
ground) is what it appears to do inside the frame \emph{plus} the frame's own acceleration:
$\va_{\text{ground}}=\va'+\va_0$.

\stepn{2} \textbf{Apply the real law in the ground frame,} where it is valid:
$\vF_{\text{real}}=m\va_{\text{ground}}=m\va'+m\va_0$.

\stepn{3} \textbf{Rearrange to look like $\vF=m\va'$ inside the frame.} Move the $m\va_0$ term to the
force side:
\[
\vF_{\text{real}}\;\underbrace{-\,m\va_0}_{\vF_{\text{pseudo}}}=m\va'.
\]
So Newton's law \emph{does} work inside the accelerating frame, provided you add one extra force
$\vF_{\text{pseudo}}=-m\va_0$ to every body, pointing \emph{opposite} to the frame's acceleration, with
magnitude $m a_0$. Add it, then solve exactly as if the frame were at rest.
\end{derivbox}

\begin{trapbox}{a pseudo-force has no source and no reaction}
A pseudo-force exists only because the frame accelerates. It has no source body, no third-law partner,
and vanishes the instant you return to an inertial frame. Choose one frame and stay in it: in the
ground frame use real forces only and write $\vF_{\text{real}}=m\va_{\text{ground}}$; in the accelerating
frame add $-m\va_0$ and use $\va'$. Writing both the pseudo-force and the frame's acceleration in one
equation counts that acceleration twice.
\end{trapbox}

\begin{figure}[htbp]
\centering
\tikzfig{pseudo_pendulum}
\figcaption{fig:pseudo}{A bob hangs in a car accelerating at $a_0$. In the car's (non-inertial) frame it
is in equilibrium under tension, weight, and a backward pseudo-force $ma_0$, giving $\tan\theta=a_0/g$.
The ground frame, with no pseudo-force, gets the same angle.}
\end{figure}

Figure~\ref{fig:pseudo} shows the simplest application. A bob hangs from the roof of a car that
accelerates horizontally at $a_0$, and the string settles at an angle $\theta$ to the vertical.
\emph{In the car's frame} the bob is at rest under the tension, the weight and a backward pseudo-force
$ma_0$: $T\sin\theta=ma_0$ and $T\cos\theta=mg$, so $\tan\theta=a_0/g$. \emph{In the ground frame}
there is no pseudo-force; the bob accelerates forward at $a_0$, so $T\sin\theta=ma_0$ is now Newton's
second law rather than a balance, and $T\cos\theta=mg$ as before. The equations, and the angle, are the
same; only their interpretation differs.

\begin{keybox}{Effective gravity}
Inside a lift accelerating up at $a_0$, every body carries an extra pseudo-force $ma_0$ downward, so it
behaves as if gravity were $g_{\text{eff}}=g+a_0$. Accelerating down, $g_{\text{eff}}=g-a_0$ (and zero in
free fall). To solve almost any in-lift problem (apparent weight, pulleys, projectiles, pendulum
periods), replace $g$ by $g_{\text{eff}}$. In general, in a frame translating with acceleration $\va_0$,
$\vg_{\text{eff}}=\vg-\va_0$: in the car of Figure~\ref{fig:pseudo} its magnitude is
$\sqrt{g^{2}+a_0^{2}}$, and the bob hangs along it.
\end{keybox}

\begin{exbox}{Apparent weight in a lift, from inside the lift}
A person of mass $m$ stands on a scale in a lift accelerating up at $a_0$. \chref{ch:newton} found the
scale reading in the ground frame: $N-mg=ma_0$. In the lift's frame the person is at rest under $N$, the
weight $mg$ and the pseudo-force $ma_0$ downward, so $N-mg-ma_0=0$. Either way $N=m(g+a_0)=mg_{\text{eff}}$:
the scale reads the effective weight. Accelerating down, $N=m(g-a_0)$, and in free fall ($a_0=g$)
$N=0$. \ans{$N=m(g\pm a_0)$}
\end{exbox}

\subsection{Working inside a rotating frame}\label{sub:rotframe}
The repair above assumed a frame that does not rotate, so that every point of it shares one acceleration
$\va_0$. The passenger and the rotor rider of this section's opening sit in frames that \emph{do} rotate.
Some circular-motion problems become much simpler from there, at the price of one extra pseudo-force.

\begin{motivbox}{When is it worth leaving the ground frame?}
The ground frame is always \emph{correct}, but it is not always \emph{convenient}. Consider a bead free
to slide along a radial groove cut in a turntable that spins at constant $\omega$. Seen from the
ground, the bead traces an outward spiral (Figure~\ref{fig:nifgroove}, left), and the only real horizontal force, the sideways push of
the groove's wall, is perpendicular to the \emph{groove} but not to the path: it does work on the bead
and changes both the size and the direction of its velocity, which makes the ground-frame algebra
unpleasant.

Now sit \emph{on} the turntable. In that frame the bead slides straight out along the groove
(Figure~\ref{fig:nifgroove}, right).
The price is one pseudo-force, the \textbf{centrifugal force} $m\omega^{2}x$ on a bead at distance $x$
from the axis, pointing radially outward. The rule extends the one for a translating frame: add $-m\va_0$,
where $\va_0$ is now the acceleration of the point of the frame at which the body currently sits. For a frame
turning at constant $\omega$ that point has centripetal acceleration $\omega^{2}x$, so the pseudo-force
is $m\omega^{2}x$ outward. A body moving relative to the frame needs one further term, the Coriolis
force (see the remark at the end of this subsection).
\end{motivbox}

\begin{trapbox}{centrifugal force and $mv^2/r$ never share an equation}
In the ground frame, set the inward resultant of the real forces equal to $mv^2/r$. In the rotating
frame, add the outward centrifugal force $m\omega^2r$ (equal in size to $mv^2/r$) to the \emph{same} real
forces and, for a body at rest in that frame, set the total to zero. Writing both counts the circular motion twice. A passenger in a car
rounding a bend feels thrown outward, but in the ground frame nothing pushes them outward: their body
tends to go straight on, and the door pushes them \emph{inward} round the bend.
\end{trapbox}

\begin{figure}[htbp]
\centering
\tikzfig{nif_groove}
\figcaption{fig:nifgroove}{The same bead described twice. In the ground frame (left) the wall's push $N$ is the only horizontal force; the bead's velocity is not along the groove, and its path is a spiral. In the table's frame (right) the table is at rest, $N$ is balanced across the groove by the Coriolis pseudo-force, and the only force along the groove is the centrifugal pseudo-force $m\omega^2x$: the motion is a straight slide along $x$.}
\end{figure}

\begin{exbox}{Bead sliding in a groove on a rotating table}
A horizontal table turns at constant angular speed $\omega$ about its axis. A smooth radial groove is
cut in it, and a particle is gently placed inside the groove at a distance $a$ from the centre. Find
the particle's speed \emph{relative to the table} when its distance from the centre has grown to $b$.

\stepn{1} \textbf{Choose the frame deliberately.} Work in the frame of the \emph{table}. The particle
then moves along a straight line (the groove), and we must include the centrifugal pseudo-force.

\stepn{2} \textbf{Write the force along the groove.} At distance $x$ from the axis the outward
centrifugal force is $m\omega^{2}x$. The groove is smooth, and its wall pushes only \emph{sideways}, so
along the groove this is the only force:
\[
ma = m\omega^{2}x
\qquad\Longrightarrow\qquad
a=\omega^{2}x.
\]

\stepn{3} \textbf{The acceleration depends on position, not time}, so use the chain-rule form
$a=\dfrac{dv}{dt}=\dfrac{dv}{dx}\dfrac{dx}{dt}=v\dfrac{dv}{dx}$ rather than any constant-acceleration
formula:
\[
v\frac{dv}{dx}=\omega^{2}x.
\]

\stepn{4} \textbf{Separate and integrate} from $(x=a,\ v=0)$ to $(x=b,\ v)$:
\[
\int_{0}^{v} v\,dv=\omega^{2}\int_{a}^{b} x\,dx
\qquad\Longrightarrow\qquad
\frac{v^{2}}{2}=\frac{\omega^{2}}{2}\left(b^{2}-a^{2}\right)
\]
\[
\boxed{\;v=\omega\sqrt{b^{2}-a^{2}}\;}
\]
\ans{$v=\omega\sqrt{b^{2}-a^{2}}$}

\smallskip
\textbf{Checks.} At $b=a$ the formula gives $v=0$ \checkmark. A particle placed exactly at the centre
would stay there, because the centrifugal force vanishes at the axis: an unstable equilibrium, but an
equilibrium.
\end{exbox}

\begin{figure}[htbp]
\centering
\tikzfig{nif_centriwork}
\figcaption{fig:nifwork}{The centrifugal force depends only on the distance from the axis and points
radially, so the work it does between two radii is the same along every path: it is a conservative
force within the rotating frame.}
\end{figure}

\begin{exbox}{Work done by the centrifugal force \textnormal{(a classic problem)}}
A small body of mass \SI{0.10}{kg} moves in a frame rotating at a constant
$\omega=\SI{5.0}{rad/s}$. What work does the centrifugal force perform as the body is carried along an
\emph{arbitrary} path from $r_1=\SI{30}{cm}$ to $r_2=\SI{50}{cm}$ from the rotation axis?

\stepn{1} \textbf{Why ``arbitrary path'' is not a problem.} The centrifugal force $F=m\omega^{2}r$ is
directed radially outward and its magnitude depends only on $r$. A force of that form is
\textbf{conservative} (a term made precise in \chref{ch:energy}): only the radial part of any
displacement does work, so all the sideways
wandering contributes nothing and only the change in $r$ matters (Figure~\ref{fig:nifwork}).

\stepn{2} \textbf{Integrate along the radius.}
\[
W=\int_{r_1}^{r_2} m\omega^{2}r\,dr
= m\omega^{2}\left[\frac{r^{2}}{2}\right]_{r_1}^{r_2}
=\frac{m\omega^{2}\left(r_2^{2}-r_1^{2}\right)}{2}.
\]

\stepn{3} \textbf{Substitute.}
\[
W=\frac{(0.10)(5.0)^{2}\big[(0.50)^{2}-(0.30)^{2}\big]}{2}
=\frac{(0.10)(25)(0.25-0.09)}{2}
=\frac{2.5\times0.16}{2}=\boxed{\SI{0.20}{J}}
\]
\ans{$W=\SI{0.20}{J}$}

\smallskip
\textbf{Positive, as it must be:} the body moved \emph{outward}, in the same direction as the
centrifugal force, so that force did positive work on it.
\end{exbox}

\begin{remarkbox}{A word on the Coriolis force}
A rotating frame requires \emph{two} pseudo-forces. The centrifugal force $m\omega^{2}r$ acts on
\emph{every} body in the frame, moving or not. A body \textbf{moving relative to the frame} needs a
second correction, the \textbf{Coriolis force} $2m(\vv_{\text{rel}}\ex\vomega)$, where $\vv_{\text{rel}}$ is the
velocity relative to the frame and $\vomega$ the frame's angular velocity vector, along the axis
(\chref{ch:circkin}); it is perpendicular to both. In the groove problem it points across the groove, so the wall's normal force
absorbs it and it never enters the along-groove equation; that is why we could ignore it. On the
rotating Earth the Coriolis force deflects winds, giving cyclones their spin, and bends the paths of
long-range projectiles. It lies beyond this syllabus.
\end{remarkbox}

\subsection{Artificial gravity: the rotor, scaled up to a spacecraft}

The ground-frame and rotating-frame descriptions now meet in one practical problem. A spacecraft coasting between planets is in \textbf{free fall}. The ship and everyone in it accelerate
together under whatever gravity is present, so nobody presses on anything, and the crew float. That is
the \emph{default} state, and it is a problem: months of it waste muscle and leach calcium from bone.
Spinning the ship restores a sensation of weight. This is why, in both \emph{Interstellar} and \emph{The Martian}, the astronauts float in the
ship's central hub and walk normally in the rotating ring: the same craft gives two different
sensations, and the only difference is the radius.


\begin{derivbox}{Why a spinning floor feels like a planet}
This is the rotor of \S\ref{sub:rotor} with gravity switched off, so that the wall becomes the floor. An astronaut
stands on the \emph{inside} of the outer rim, feet outward, head towards the axis.

\stepn{1} \textbf{Apply the central principle.} The astronaut travels a circle of radius $r$ at angular
speed $\omega$, so some real force must supply $m\omega^{2}r$ inward. The only thing in contact with
them is the floor, so the floor must provide all of it:
\[
N=m\omega^{2}r \qquad\text{directed \emph{inward}, towards the axis.}
\]

\stepn{2} \textbf{Invoke Newton's third law.} The astronaut therefore pushes \emph{outward} on the floor
with the same magnitude $m\omega^{2}r$ (Figure~\ref{fig:artgrav}, left).

\stepn{3} \textbf{Compare with standing on Earth.} There the floor pushes up on you with $N=mg$ and you
press down on it with $mg$. At the point of contact the two situations are mechanically identical: a
surface pushes on your feet and your skeleton carries that push up through your body. The sensation we
\emph{call} weight is that compression, so, standing still, your body cannot tell the two apart.
(Moving, or comparing head with feet, it can, as the remark below shows.) Setting $m\omega^{2}r=mg_{\text{eff}}$,
the effective gravity is
\[
\boxed{\;g_{\text{eff}}=\omega^{2}r=\frac{v^{2}}{r}\;}
\]
\end{derivbox}

\why{\textbf{Notice what is absent from that derivation: any outward force on the astronaut.} The only
real force is the floor's \emph{inward} push, which stops the crew moving off in a straight line. In
the rotating frame of \S\ref{sub:rotframe} you would add the centrifugal force $m\omega^{2}r$ outward
and solve a statics problem: the same answer, but only the ground frame identifies a real cause.}

\begin{figure}[htbp]
\centering
\tikzfig{nif_artgrav}
\figcaption{fig:artgrav}{\textbf{Left:} the mechanism. The floor of the rim pushes the astronaut inward
with $N=m\omega^{2}r$, and by the third law the astronaut presses outward on it; that mutual push is
what the body registers as weight. At the hub, where $r\approx0$, there is nothing to stand on and the
crew float. \textbf{Right:} the design penalty. Since $g_{\text{eff}}=\omega^{2}r$, a given gravity can
be bought with a small fast ring or a large slow one, but comfort forces designers towards the latter.}
\end{figure}

\begin{exbox}{The \emph{Endurance} from \emph{Interstellar}}
The ring of the \emph{Endurance} is \SI{64}{m} across and turns at \SI{5.5}{rpm}. What gravity does its
crew experience, and what does a \SI{1.8}{m} astronaut feel between head and feet? ($g=\SI{9.8}{m/s^2}$.)

\stepn{1} \textbf{Convert the spin rate to \si{rad/s}.}
\[
\omega=\frac{2\pi n}{60}=\frac{2\pi(5.5)}{60}=\SI{0.576}{rad/s}.
\]

\stepn{2} \textbf{Effective gravity at the rim} ($r=\SI{32}{m}$):
\[
g_{\text{eff}}=\omega^{2}r=(0.576)^{2}(32)=\SI{10.6}{m/s^2}\approx\mathbf{1.08\,g}.
\]

\stepn{3} \textbf{The head-to-foot gradient.} Because $g_{\text{eff}}\propto r$ and the head, a height
$h=\SI{1.8}{m}$ above the feet, is nearer the axis,
\[
\frac{g_{\text{head}}}{g_{\text{feet}}}=\frac{r-h}{r}=\frac{32-1.8}{32}=0.944,
\]
so the head feels about $\mathbf{5.6\%}$ less gravity than the feet. On Earth the corresponding figure
is about six parts in ten million. \ans{$1.08\,g$; a $5.6\%$ head-to-foot gradient}
\end{exbox}

\begin{remarkbox}{Why real designs have to be enormous}
Since $g_{\text{eff}}=\omega^{2}r$, the same gravity can be had from a small ring spun fast or a large
ring spun slowly. Three effects punish the fast option, and all three grow as $\omega$ grows.

\medskip
\textbf{1. Coriolis force.} Anything moving relative to the rotating frame feels
$2m(\vv_{\text{rel}}\ex\vomega)$, perpendicular to both its velocity and the axis. An astronaut
climbing a ladder towards the hub at \SI{1.5}{m/s} inside the \emph{Endurance} is pushed sideways,
along the rim, with
\[
a_{\text{cor}}=2\omega v_{\text{rel}}=2(0.576)(1.5)=\SI{1.73}{m/s^2},
\]
about $18\%$ of $g$. Walking along the rim instead, the same force points radially and makes you
heavier or lighter: that is effect~3 below, seen from the rotating frame. Worse, every turn of the head
moves the fluid in the inner
ear in a misleading way, which the brain reads as vertigo and nausea. Studies put the comfortable
ceiling near \SI{2}{rpm}, with \SIrange{4}{6}{rpm} tolerable only after adaptation.

\medskip
\textbf{2. Gravity gradient.} As computed above, $g_{\text{eff}}\propto r$ makes your head lighter than
your feet by the fraction $h/r$. That is $5.6\%$ on a \SI{32}{m} ring and would be $18\%$ on a
\SI{10}{m} one: enough to feel on every stand-up.

\medskip
\textbf{3. Walking changes your weight.} Move at speed $u$ along the rim and your true speed becomes
$\omega r\pm u$, so
\[
g_{\text{eff}}=\frac{(\omega r\pm u)^{2}}{r}\approx\omega^{2}r\left(1\pm\frac{2u}{\omega r}\right),
\]
by the binomial approximation of \chref{ch:toolkit}, since $u\ll\omega r$.
On the \emph{Endurance} (rim speed $\omega r\approx\SI{18.4}{m/s}$), walking at \SI{1.5}{m/s} with
the spin makes you about $17\%$ heavier and walking against it about $16\%$ lighter. Jogging one way
round the ring would be noticeably harder than the other.

\medskip
\begin{center}
\begin{tabular}{@{}lll@{}}
\toprule
Spin rate & Radius for $1\,g$ & Rim speed\\
\midrule
\SI{1}{rpm}   & \SI{894}{m} & \SI{94}{m/s}\\
\SI{2}{rpm}   & \SI{223}{m} & \SI{47}{m/s}\\
\SI{3}{rpm}   & \SI{99}{m}  & \SI{31}{m/s}\\
\SI{5.5}{rpm} & \SI{30}{m}  & \SI{17}{m/s}\\
\bottomrule
\end{tabular}
\end{center}

\medskip
This is why serious station concepts are hundreds of metres across. To deliver $1\,g$ at a
comfortable \SI{2}{rpm}, the \emph{Endurance} would need to be roughly \textbf{seven times} its filmed
diameter.
\end{remarkbox}

\begin{trapbox}{the ISS is weightless for a quite different reason}
Astronauts float on the International Space Station not because it is far from Earth, where gravity is
still about $90\%$ of its surface value, and not because it spins. They float because the station is in
\textbf{free fall}: it and everything in it accelerate together, so nothing presses on anything. The ISS
has no rotating section. Attach a ring of radius \SI{32}{m} spinning at \SI{5.5}{rpm} and its crew
would stand on its floor at once.
\end{trapbox}

\section{Summary}\label{sec:summary10}
In every circular-motion problem the inward net force, made of real forces only, equals $mv^2/r$; if the
speed also changes, the tangential net force equals $ma_t$. The settings differ only in which real force
does the job. In an accelerating frame, and only there, a pseudo-force is added. The table collects the
results.

\medskip
{\small\renewcommand{\arraystretch}{1.45}
\begin{center}
\begin{tabular}{@{}ll@{}}
\toprule
Situation & Key result\\
\midrule
Conical pendulum      & $v=\sqrt{rg\tan\theta}$, \; $T_{\text{per}}=2\pi\sqrt{\dfrac{\ell\cos\theta}{g}}$, \; $r=\ell\sin\theta$\\[3pt]
Inside a smooth cone  & $v=\sqrt{rg\tan\theta}$, $\theta$ = wall's angle to the horizontal\\[3pt]
Leaning cyclist       & $\tan\theta=\dfrac{v^{2}}{rg}$\\[3pt]
Rotor / death well    & $\omega_{\min}=\sqrt{\dfrac{g}{\mu r}}$ \; (friction vertical, $N$ horizontal)\\[3pt]
Flat road             & $v_{\max}=\sqrt{\mu rg}$\\[3pt]
Banked, smooth        & $\tan\theta=\dfrac{v^{2}}{rg}$: one design speed only\\[3pt]
Banked with friction  & $v_{\max}=\sqrt{rg\dfrac{\tan\theta+\mu}{1-\mu\tan\theta}}$, \;
                        $v_{\min}=\sqrt{rg\dfrac{\tan\theta-\mu}{1+\mu\tan\theta}}$\\[5pt]
Toppling              & $v_{\text{topple}}=\sqrt{\dfrac{gra}{h}}$ \;: the \emph{second} critical speed\\[4pt]
Vertical circle       & $v^{2}=u^{2}-2gh$ \;and\; $T=\dfrac{mu^{2}}{R}+mg(3\cos\theta-2)$\\[4pt]
\quad complete        & $u\ge\sqrt{5gR}$ \; (string) \quad or \quad $u>2\sqrt{gR}$ \; (rod/tube)\\[2pt]
\quad leaves          & $\sqrt{2gR}<u<\sqrt{5gR}$, slack at $h_2=\dfrac{u^{2}+gR}{3g}$\\[3pt]
\quad oscillates      & $u\le\sqrt{2gR}$, stops at $h_1=\dfrac{u^{2}}{2g}$\\[3pt]
\quad always          & $T_b-T_t=6mg$\\[3pt]
Ball off a sphere     & leaves at $\cos\theta=\dfrac23+\dfrac{u^{2}}{3gR}$; from rest, $48.2^\circ$\\[3pt]
Accelerating frame    & add $-m\va_0$ to every body; in a lift $g_{\text{eff}}=g\pm a_0$\\[3pt]
Rotating frame        & add $m\omega^{2}r$ outward; $W=\dfrac{m\omega^{2}(r_2^{2}-r_1^{2})}{2}$\\[3pt]
Artificial gravity    & $g_{\text{eff}}=\omega^{2}r$\\
\bottomrule
\end{tabular}
\end{center}
}

\medskip
\textbf{Three questions that settle almost everything.}
\emph{(1)} Which real force is doing the turning?
\emph{(2)} Can the constraint push, or only pull? (String pulls only; rod, tube and groove do both.)
\emph{(3)} If friction vanished, which way would the body slide? Friction points the other way.

\section{Problem Types}\label{sec:types10}
Circular dynamics questions are dominated by two set pieces: the banked road and the vertical circle.
The five types below cover those and the closely related conical pendulum and rotor, in the order of the
chapter. Every one is solved
by the same opening: draw the real forces, choose horizontal (radial) and vertical axes, which put the
whole acceleration along one axis, and set the inward resultant equal to $mv^2/r$.

\archq{Conical pendulum}
{``A bob whirls in a horizontal circle on a string making angle $\theta$ with the vertical, find the
tension, speed or period.''}
{Resolve: vertical $T\cos\theta=mg$; horizontal $T\sin\theta=\dfrac{mv^2}{r}$ with $r=\ell\sin\theta$.
Dividing gives $\tan\theta=\dfrac{v^2}{rg}$, and $T=\dfrac{mg}{\cos\theta}$, always \emph{greater}
than the weight.}
{Setting $T=mg$. The string must both support the bob \emph{and} supply the centripetal pull, so it is
always tighter than the weight.}
{A conical pendulum's string makes $60^\circ$ with the vertical. The tension is:}
{$mg/2$}{$mg$}{$2mg$}{$\sqrt3\,mg$\archans{C}{$T=\dfrac{mg}{\cos60^\circ}=\dfrac{mg}{0.5}=2mg$. (A) multiplies
by $\cos60^\circ$ instead of dividing; (B) is the trap $T=mg$.}}

\archq{The rotor and death well}
{``A person is pressed against the inside wall of a rotating cylinder, find the minimum angular speed
so they do not slide down.''}
{The \emph{normal} force points inward and supplies the centripetal force, $N=m\omega^2r$; \emph{friction}
acts \emph{upward} and must support the weight, $\mu N\ge mg$. Combining,
$\boxed{\omega_{\min}=\sqrt{\dfrac{g}{\mu r}}}$, independent of mass.}
{Taking friction as horizontal. Here the wall is vertical, so the normal force is horizontal and
friction is \emph{vertical}: the reverse of the usual floor case.}
{In a rotor of radius \SI{2}{m} with $\mu=0.5$ ($g=\SI{10}{m/s^2}$), the minimum angular speed is:}
{\SI{1}{rad/s}}{\SI{2.2}{rad/s}}{\SI{3.2}{rad/s}}{\SI{10}{rad/s}\archans{C}{
$\omega_{\min}=\sqrt{\dfrac{g}{\mu r}}=\sqrt{\dfrac{10}{0.5\times2}}=\sqrt{10}\approx\SI{3.2}{rad/s}$.
(B) leaves out $\mu$; (D) forgets the square root.}}

\archq{Banked or level road}
{``Find the angle of banking for a road of radius $r$ at speed $v$,'' or the maximum safe speed on a
banked or level road.}
{For a frictionless banked road the horizontal component of the normal force supplies the centripetal
force: $N\sin\theta=\dfrac{mv^2}{r}$ and $N\cos\theta=mg$. Dividing,
$\boxed{\tan\theta=\dfrac{v^2}{rg}}$, independent of mass. On a \emph{level} road friction must supply it
instead, giving $v_{\max}=\sqrt{\mu rg}$.}
{Writing $N=mg\cos\theta$ as on an inclined plane. The acceleration is \emph{horizontal}, so it has a
component $\dfrac{v^2}{r}\sin\theta$ perpendicular to the slope, and
$N=mg\cos\theta+\dfrac{mv^2}{r}\sin\theta$.}
{The angle of banking for a curve of radius \SI{100}{m} designed for \SI{20}{m/s} is given by
($g=\SI{10}{m/s^2}$):}
{$\tan\theta=0.2$}{$\tan\theta=0.4$}{$\tan\theta=2$}{$\tan\theta=4$\archans{B}{
$\tan\theta=\dfrac{v^2}{rg}=\dfrac{400}{100\times10}=0.4$, that is, $\theta\approx21.8^\circ$.
(A) is $v/r$, which is not even dimensionless; (D) forgets $g$.}}

\archq{Banked road with friction}
{``Find the maximum (or minimum) safe speed on a banked road when friction is also present.''}
{With friction acting \emph{down} the slope at the limit,
$v_{\max}=\sqrt{rg\dfrac{\tan\theta+\mu}{1-\mu\tan\theta}}$; with friction up the slope,
$v_{\min}=\sqrt{rg\dfrac{\tan\theta-\mu}{1+\mu\tan\theta}}$. Setting $\mu=0$ recovers the frictionless
$v=\sqrt{rg\tan\theta}$.}
{Using the frictionless formula when $\mu$ is given, or getting the sign of the friction term
backwards. At maximum speed the car tends to slide \emph{outward and up}, so friction acts \emph{down}
the slope.}
{A curve of radius \SI{90}{m} is banked with $\tan\theta=0.5$, and $\mu=1/3$ ($g=\SI{10}{m/s^2}$). The
maximum safe speed is:}
{\SI{11.3}{m/s}}{\SI{21.2}{m/s}}{\SI{27.4}{m/s}}{\SI{30}{m/s}\archans{D}{
$\dfrac{\tan\theta+\mu}{1-\mu\tan\theta}=\dfrac{5/6}{5/6}=1$, so $v_{\max}=\sqrt{rg}=\sqrt{900}=\SI{30}{m/s}$.
(A) is $v_{\min}$, (B) the frictionless design speed, (C) the result of forgetting the denominator.}}

\archq{Vertical circle, critical speeds and tensions}
{``A body on a string/track goes round a vertical circle, find the minimum speed at the top or
bottom, or the tension somewhere.''}
{Two laws together: Newton along the radius \emph{and} energy conservation between the two points. At
the top both $T$ and $mg$ point inward, so $T+mg=\dfrac{mv_t^2}{R}$; setting $T\ge0$ gives
$v_t^{\min}=\sqrt{gR}$, and energy then gives $v_b^{\min}=\sqrt{5gR}$. Always
$T_b-T_t=6mg$.}
{Using $v_{\min}=\sqrt{gR}$ at the \emph{bottom}. That is the \emph{top} condition. Also, for a rigid
\emph{rod} (which can push) the top condition relaxes to $v_t>0$, so the speed at the bottom need only
exceed $2\sqrt{gR}$.}
{A stone on a string of length \SI{1.6}{m} is whirled in a vertical circle ($g=\SI{10}{m/s^2}$). The minimum speed at
the highest point is:}
{\SI{2}{m/s}}{\SI{4}{m/s}}{\SI{8}{m/s}}{\SI{16}{m/s}\archans{B}{$v_t^{\min}=\sqrt{gR}
=\sqrt{10\times1.6}=\sqrt{16}=\SI{4}{m/s}$. (C) is $2\sqrt{gR}$, the bottom speed for a \emph{rod};
(D) forgets the square root.}}

\section{Practice Problems}\label{sec:practice10}

The problems below are of standard examination difficulty and are meant to be worked without looking back at the chapter. Answers, with brief working, follow the list.

\begin{enumerate}[itemsep=3pt]
\item A car rounds a \emph{level} curve of radius \SI{60}{m} with $\mu=0.6$. Find its maximum speed
      ($g=\SI{10}{m/s^2}$).
\item A curve of radius \SI{90}{m} is banked at $45^\circ$ and is frictionless. Find the design speed ($g=\SI{10}{m/s^2}$).
\item A stone on a string of length \SI{0.9}{m} is whirled in a vertical circle ($g=\SI{10}{m/s^2}$). Find (a) the
      minimum speed at the lowest point to complete the circle and (b) the corresponding speed at the
      top.
\item A conical pendulum of length \SI{1.0}{m} makes $45^\circ$ with the vertical ($g=\SI{9.8}{m/s^2}$).
      Find its period and the tension in terms of $mg$.
\item In a rotor of radius \SI{2.5}{m} with $\mu=0.5$, find the minimum angular speed for a rider not to
      slide down ($g=\SI{10}{m/s^2}$).
\item A loaded van of track width $2a=\SI{1.5}{m}$ has its centre of gravity \SI{1.5}{m} above the road and
      rounds a curve of radius \SI{50}{m} ($g=\SI{10}{m/s^2}$). Find the speed at which it would overturn, and say
      whether skidding or toppling occurs first if $\mu=0.8$.
\end{enumerate}

\subsection*{Answers}
{\small
\begin{enumerate}[itemsep=2pt]
\item $v_{\max}=\sqrt{\mu rg}=\sqrt{0.6(60)(10)}=\sqrt{360}\approx\SI{19.0}{m/s}$.
\item $v=\sqrt{rg\tan45^\circ}=\sqrt{90(10)}=\SI{30}{m/s}$.
\item (a) $u_{\min}=\sqrt{5gR}=\sqrt{5(10)(0.9)}=\sqrt{45}\approx\SI{6.7}{m/s}$;
(b) $v_{\text{top}}=\sqrt{gR}=\sqrt{9}=\SI{3}{m/s}$.
\item $T_{\text{per}}=2\pi\sqrt{\dfrac{\ell\cos\theta}{g}}=2\pi\sqrt{\dfrac{1.0(0.707)}{9.8}}
\approx\SI{1.69}{s}$; tension $=\dfrac{mg}{\cos45^\circ}=\sqrt2\,mg\approx\mathbf{1.41\,mg}$.
\item $\omega_{\min}=\sqrt{\dfrac{g}{\mu r}}=\sqrt{\dfrac{10}{0.5(2.5)}}=\sqrt{8}
\approx\SI{2.83}{rad/s}$.
\item $a=\SI{0.75}{m}$, so $v_{\text{topple}}=\sqrt{\dfrac{gra}{h}}
=\sqrt{\dfrac{10(50)(0.75)}{1.5}}=\sqrt{250}\approx\SI{15.8}{m/s}$. Meanwhile
$v_{\text{skid}}=\sqrt{\mu rg}=\sqrt{0.8(50)(10)}=\SI{20}{m/s}$. Since
$v_{\text{topple}}<v_{\text{skid}}$, this tall, narrow vehicle \textbf{overturns before it skids}, the
dangerous case ($\mu=0.8>a/h=0.5$).
\end{enumerate}
}


\section{Challenge Problems}\label{sec:challenge10}
These problems extend the methods of \chsref{ch:friction}{ch:circdyn} beyond the standard examination
question. Each states which idea it builds on. Attempt each before reading the solution.

\begin{challengebox}{The most efficient pull, and the most efficient hold (builds on optimising against friction)}
A \SI{10}{kg} crate is placed on a rough incline of $30^\circ$ ($\mu_s=0.3$). (a) Find the least force
that will drag it up the slope, and the angle at which to apply it. (b) Left alone, the crate slides
down. Find the least force that will hold it at rest, and its direction. ($g=\SI{9.8}{m/s^2}$.)
\challengeans{Let the pull make an angle $\beta$ above the incline. Perpendicular to the incline,
$N=mg\cos\theta-F\sin\beta$; along it, on the point of slipping, $F\cos\beta=mg\sin\theta+\mu N$. Hence
\[
F=\frac{mg(\sin\theta+\mu\cos\theta)}{\cos\beta+\mu\sin\beta}.
\]
With the angle of friction $\lambda=\arctan\mu=\arctan(0.3)=16.7^\circ$, the denominator is
$\cos(\beta-\lambda)/\cos\lambda$, largest at $\beta=\lambda$. So the least force is applied at
$\beta=\lambda$ \emph{above the incline surface}, and
\[
F_{\min}=mg\sin(\theta+\lambda)=10\times9.8\times\sin(46.7^\circ)\approx\SI{71.3}{N}.
\]
Pulling at the friction angle rather than straight along the slope saves force by lifting some weight off
the surface, which cuts the normal force and with it the friction. (Check: $F\sin\lambda\approx\SI{20}{N}
<mg\cos\theta\approx\SI{85}{N}$, so the crate stays on the surface.)

(b) The crate slides if left alone because $\tan30^\circ=0.577>\mu_s$. On the point of sliding
\emph{down}, friction acts \emph{up} the slope, so the signs of the friction terms reverse:
$F\cos\beta+\mu N=mg\sin\theta$ with $N=mg\cos\theta-F\sin\beta$, giving
\[
F=\frac{mg(\sin\theta-\mu\cos\theta)}{\cos\beta-\mu\sin\beta}.
\]
The denominator is now largest at $\beta=-\lambda$: the force should be directed at the friction angle
\emph{into} the slope, pressing the crate on to the surface to increase the friction that helps. Then
$F_{\min}=mg\sin(\theta-\lambda)=98\sin(13.3^\circ)\approx\SI{22.5}{N}$. Lifting helps when
friction opposes you; pressing helps when friction is on your side.}
\end{challengebox}

\begin{challengebox}{The full speed band of a banked curve (builds on banking with friction)}
A curve of radius \SI{50}{m} is banked at $15^\circ$, with $\mu_s=0.5$ between tyres and road
($g=\SI{9.8}{m/s^2}$). (a) Find the maximum speed at which a car can round it without sliding.
(b) Show that there is no minimum speed. (c) Keeping $\mu_s=0.5$, above what bank angle would there be
no maximum speed either?
\challengeans{(a) With friction acting down-slope at maximum speed,
\[
v_{\max}=\sqrt{rg\,\frac{\mu_s+\tan\theta}{1-\mu_s\tan\theta}}
=\sqrt{50\times9.8\times\frac{0.5+\tan15^\circ}{1-0.5\tan15^\circ}}\approx\SI{20.8}{m/s}\approx\SI{75}{km/h}.
\]
Both effects help at the top end: the banking contributes an inward $N\sin\theta$, and friction adds its
share on top. Without the bank ($\theta=0$) the limit would be only $\sqrt{\mu_s rg}=\SI{15.7}{m/s}$;
the $15^\circ$ tilt buys nearly \SI{20}{km/h}.

(b) Here $\mu_s=0.5>\tan15^\circ=0.268$, so $\tan\theta-\mu_s<0$ and the formula for $v_{\min}^{2}$
is negative: friction alone can hold a stationary car on this bank, and no speed is too slow.

(c) $v_{\max}$ becomes unbounded when the denominator $1-\mu_s\tan\theta$ reaches zero, at
$\tan\theta=1/\mu_s=2$, that is $\theta\approx63.4^\circ$. On a steeper bank the vertical balance
$N(\cos\theta-\mu_s\sin\theta)=mg$ can be met by no positive $N$ when friction acts down the slope at
its limit: friction can always grow enough to stop an upward slide. (A stationary car would slide down such a
bank, since $\tan\theta>\mu_s$.) The extreme case $\theta=90^\circ$ is the death well of
\S\ref{sub:rotor}, where only a minimum speed exists.}
\end{challengebox}

\begin{challengebox}{A string that can break (builds on vertical circle)}
A \SI{0.5}{kg} ball on a string of length \SI{2}{m} is whirled in a vertical circle. The string breaks if
its tension exceeds \SI{50}{N}. (a) Find the range of speeds $u$ at the lowest point for which the ball goes
all the way round with the string neither slackening nor breaking. (b) If instead $u=\SI{8.0}{m/s}$,
find the angle from the lowest point at which the string goes slack. ($g=\SI{9.8}{m/s^2}$.)
\challengeans{(a) Two conditions, one at each end of the circle. The tension is \emph{least} at the top, so
slackening is tested there: $T_{\text{top}}=\dfrac{mu^{2}}{R}-5mg\ge0$ gives
$u\ge\sqrt{5gR}=\sqrt{98}=\SI{9.90}{m/s}$. The tension is \emph{greatest} at the bottom, so breaking is
tested there:
\[
T_{\text{bottom}}=\frac{mu^{2}}{R}+mg\le50
\;\Longrightarrow\;
u^{2}\le\frac{(50-4.9)(2)}{0.5}=180.4
\;\Longrightarrow\;
u\le\SI{13.4}{m/s}.
\]
So $\SI{9.90}{m/s}\le u\le\SI{13.4}{m/s}$. Because $T_{\text{bottom}}-T_{\text{top}}=6mg=\SI{29.4}{N}$
always, the window exists only because the breaking tension exceeds $6mg$: a string rated below
\SI{29.4}{N} could never take this ball round.

(b) First decide the regime: $\sqrt{2gR}=\SI{6.26}{m/s}<8.0<\SI{9.90}{m/s}$, so the string slackens
before the ball stops. From \S\ref{sub:vcfail}, $h_2=\dfrac{u^{2}+gR}{3g}=\dfrac{64+19.6}{29.4}
=\SI{2.84}{m}$, so $\cos\theta=\dfrac{R-h_2}{R}=-0.422$ and $\theta\approx115^\circ$, in the upper
half as the theory requires.}
\end{challengebox}

\section*{Looking Back on Part B}
\addcontentsline{toc}{section}{Looking Back on Part B}

\textbf{Part B gave motion its cause.}
\begin{itemize}
\item \textbf{Newton's laws (\chref{ch:newton}).} The first law defines force and the inertial frame; the second,
$\vF=m\va$ (more generally $\vF=d\vp/dt$), does almost all the work, component by component; the third pairs every
force with an equal, opposite one on a \emph{different} body, which is why action--reaction pairs never
cancel and why propulsion is possible. The free-body diagram, every real force on one body, resolved to
axes, is the master technique.
\item \textbf{Connected bodies (\chref{ch:connected}).} An ideal string carries one tension throughout and links its
bodies by an acceleration constraint. Atwood, block-on-table, and incline-pulley are one machine at three
tilts. Springs combine by asking what is \emph{shared}: same force $\Rightarrow$ series (softer), same
extension $\Rightarrow$ parallel (stiffer).
\item \textbf{Friction (\chref{ch:friction})} is self-adjusting up to a ceiling $\mu_s N$, then drops to $\mu_k N$; the
angle of repose measures $\mu_s$ directly, and the least force to drag a body is applied at the angle of
friction.
\item \textbf{Circular dynamics (\chref{ch:circdyn}).} Centripetal force is never new: it is the net inward
component of real forces, whether friction on a flat road, the normal force on a banked one, or tension
in a vertical circle, where energy conservation must supply the speed. In an accelerating frame, add one
pseudo-force $-m\va_0$ (in a rotating frame, the centrifugal force) and Newton's laws work again.
\end{itemize}
\textbf{The thread into Part C:} solving $\vF=m\va$ through every instant is powerful but laborious. Next
we find shortcuts, conserved quantities that connect a motion's endpoints without the middle.
